%% ======================================================================
%%  REVTeX 4.2 template -- Physical Review A / X / X Quantum
%%      pra          -> Physical Review A
%%      prx          -> Physical Review X
%%      prxquantum   -> PRX Quantum
%% ======================================================================
\documentclass[aps,prx,reprint,floatfix,amsmath,amssymb,longbibliography]{revtex4-2}
%\documentclass[aps,pra,reprint,amsmath,amssymb,longbibliography]{revtex4-2}
%\documentclass[aps,prxquantum,reprint,amsmath,amssymb,longbibliography]{revtex4-2}

%% Package Import %%%%%%%%%%%%%%%%%%%%

\usepackage{natbib}
\usepackage{listings}
\usepackage{color}
\definecolor{dkgreen}{rgb}{0,0.6,0}
\definecolor{gray}{rgb}{0.5,0.5,0.5}
\definecolor{mauve}{rgb}{0.58,0,0.82}

\lstset{frame=tb,
  language=Python,
  aboveskip=3mm,
  belowskip=3mm,
  showstringspaces=false,
  columns=flexible,
  basicstyle={\small\ttfamily},
  numbers=none,
  numberstyle=\tiny\color{gray},
  keywordstyle=\color{blue},
  commentstyle=\color{dkgreen},
  stringstyle=\color{mauve},
  breaklines=true,
  breakatwhitespace=true,
  tabsize=3
}

\usepackage{amsmath,amssymb,mathrsfs}
\usepackage{latexsym}
\usepackage{grffile}
\usepackage{graphicx,epstopdf,color}
\usepackage{amsfonts}
\usepackage[hypertexnames=false]{hyperref}
\usepackage[dvipsnames]{xcolor}
\usepackage{url}
\usepackage{soul}
\usepackage{cancel}
\usepackage{subcaption}
\usepackage{bm}
\usepackage{xr}
\usepackage{array}
\usepackage{mathtools}
\usepackage{booktabs}
\usepackage{cleveref}
\usepackage[american]{circuitikz}
\usetikzlibrary{calc}

\definecolor{bufcol}{RGB}{245,95,92}
\definecolor{cqcol}{RGB}{20,170,235}
\definecolor{gcol}{RGB}{130,130,130}

% --- multi-buffer circuit (fig_multibuffer) -------------------------------
\tikzset{
  ats/.style={draw=bufcol, thick, fill=bufcol!10, text=bufcol,
              minimum width=1.1cm, minimum height=0.75cm, inner sep=1pt,
              font=\small},
  port/.style={draw=bufcol, fill=gray!30, text=bufcol,
               minimum width=1.0cm, minimum height=0.6cm, inner sep=2pt,
               font=\small},
}
% \buffercell{x}{n}: buffer n with its ATS at abscissa x
% (memory bus at y = 4, buffer node at y = 2.5, ground rail at y = 0)
\newcommand{\buffercell}[2]{%
  \draw (#1,4) to[C, l_={$C_{c,#2}$}, *-] (#1,2.5);
  \begin{scope}[color=bufcol]
    \fill (#1,2.5) circle (1.5pt);
    \draw (#1,2.5) -- (#1,1.85);
    \node[ats] (A#2) at (#1,1.45) {ATS$_{#2}$};
    \draw (A#2.south) -- (#1,0);
    \draw (#1,2.5) -- (#1-1,2.5) to[C, l_={$C_{b,#2}$}, color=bufcol] (#1-1,0);
    \draw (#1,2.5) to[C, l={$C_{p,#2}$}, color=bufcol] (#1+1.3,2.5);
    \node[port, anchor=west] at (#1+1.3,2.5) {$\kappa_{b_{#2}},\,\eps_{d_{#2}}$};
    \node[font=\scriptsize, anchor=west] at (A#2.east) {$\eps_{#2}(t)$};
    \node[font=\small] at (#1-0.5,-0.4) {buffer $b_{#2}$};
  \end{scope}
}

%% Cleverref Settings %%%%%%%%%%%%%%%%%%%%
\crefname{equation}{Eq.}{Eqs.}
\Crefname{equation}{Equation}{Equations}
\crefname{figure}{Fig.}{Figs.}
\Crefname{figure}{Figure}{Figures}
\crefname{section}{Sec.}{Secs.}
\Crefname{section}{Section}{Sections}
\crefname{appendix}{Appendix}{Apps.}
\Crefname{appendix}{Appendix}{Apps.}
\crefname{table}{Table}{Tables}

%% New Commands %%%%%%%%%%%%%%%%%%%%%%%%
\newcommand{\textalert}[1]{}
\newcommand{\anticomm}[2]{\left\{#1,#2\right\}}
\newcommand{\vac}{\left|\,0\,\right\rangle}
\newcommand{\ket}[1]{\left|#1\right\rangle}
\newcommand{\bra}[1]{\left\langle#1\right|}
\newcommand{\bk}[1]{\bigl\langle#1\bigr\rangle}
\newcommand{\braket}[2]{\bigl\langle#1\bigl|\bigr.#2\bigr\rangle}
\newcommand{\nn}{\nonumber}
\newcommand{\pa}{\partial}
\newcommand{\up}{\uparrow}
\newcommand{\dw}{\downarrow}
\newcommand{\eps}{\varepsilon}
\newcommand{\pd}{{\phantom{\dag}}}
\newcommand{\Real}[1]{\text{Re}\left(#1\right)}
\newcommand{\Imag}[1]{\text{Im}\left(#1\right)}
\newcommand{\pb}[2]{\{#1,#2\}}
\newtheorem{theorem}{Theorem}
\newtheorem{proposition}[theorem]{Proposition}
\newcommand{\D}{\overline{\mathbb{D}}}
\newcommand{\C}{\mathbb{C}}
\newcommand{\U}{\mathcal{U}}
\newcommand{\Lc}{\mathcal{L}}

\graphicspath{{nb_fig/}}
\captionsetup[subfigure]{labelformat=brace}

\newcommand{\om}{\omega}
\newcommand{\da}{^\dagger}
\newcommand{\ha}{\hat{a}}
\newcommand{\hb}{\hat{b}}
\newcommand{\hc}{\hat{c}}

\newcommand{\norm}[1]{\lVert{#1}\rVert}

\DeclareMathOperator{\Tr}{Tr}
\renewcommand{\Im}{\text{Im}}
\renewcommand{\Re}{\text{Re}}

\newcommand{\AP}[1]{{\color{RedOrange} [AP: {#1}]}}
\newcommand{\JC}[1]{{\color{YellowOrange}[JC: {#1}]}}
\newcommand{\TMD}[1]{{\color{Blue} [TMD: {#1}]}}

\makeatletter
\def\@fnsymbol#1{\ensuremath{\ifcase#1\or * \or \mathsection\or \mathparagraph\or \|\or **\or \else\@ctrerr\fi}}
\makeatother

%%%%%%%%%%%%%%%%%%%%%%%%%%%%%%%%%%%%%%%%

\begin{document}

\title{Onset of Chaos in dissipative cat qubit}

\author{Théo Malas-Danzé}
\affiliation{Inria, ENS-PSL, Mines Paris, Alice\& Bob}
\email[Corresponding author: ]{theo.malas-danze@alice-bob.com}

\author{Joachim Cohen}
\affiliation{Alice \& Bob}

\author{Alexandru Petrescu}
\affiliation{Inria, ENS-PSL, Mines Paris}

\date{\today}

\begin{abstract}
Chaos in quantum computing is seldom a feature that is looked for: it is much harder to stabilize any state and you loose all memory of previous computation step. This is why it is crucial to investigate its onset for any architecture, and if chaos there is, how could we mitigate it. In this paper we focus on the autoparametric dissipative cat qubit, show how it can be driven into chaos, and how to partially avoid it. We set an upper threshold on the two-photon rate and buffer amplitude (cat size) and display ways to push this threshold away. Along the way, we showcase a pipeline to study chaos; the tools presented here are versatile and can be used for any Floquet systems.
\end{abstract}

\maketitle

\section{Introduction}
In order to build a useful quantum computer, one need a system which is at the same time stable, to be able to store quantum states, and non-linear, so as to perform non trivial operations. Unfortunately non-linear systems have the tendency to develop chaotic behaviors when their parameters are tuned. First investigations on chaos have been carried in chains of transmons \cite{Berke2022May}. A few years later in single qubit system such as transmons \cite{Cohen2023Apr} and Kerr cat qubit \cite{ChavezCarlos2024Nov} were demonstrated to host chaotic dynamics. A Kerr cat is made out of a single mode and quartic non-linearity and thus it is logical to expect dissipative cat qubits, which are made out of two-modes and a transcendental non-linearity to be able to develop chaos. \TMD{Cite Burgelman}

This paper not only show that chaos can develop in a specific dissipative cat qubit, but also showcase a general method to study the onset of chaos in periodic open quantum systems.

In the first part, we present the dissipative cat qubit system that we investigate. In the second part, we carry a classical analysis of this system and in the third part, we focus on the quantum analysis. Finally, we tackle two ways to mitigate chaos in dissipative cat qubit.

\section{Model / Setup}
\label{sec:model}
Dissipative cat qubits is a transatlantic superconducting architecture which has been developed in the late 2010s at Yale University \cite{Mirrahimi2014Apr} and ENS Paris \cite{Lescanne2019Jul}.

The goal of the dissipative qubit is to stabilize two coherent states $\ket{+\alpha}$ and $\ket{-\alpha}$ in quantum phase space. This is achieved through the two-photon loss Lindbladian
\begin{equation}
    \mathcal{L} \propto \mathcal{D}[\hat a^2-\alpha^2].
\end{equation}
With this Lindbladian, one of the two errors in quantum computation can be suppressed: when the cat-size $\alpha$ is increased, the bit-flip rate decreases exponentially, while the phase-flip rate only increases linearly \cite{Lescanne2019Jul}.
Such Lindbladian can be effectively engineered if you add a second mode. The mode $a$ stabilizing the cat state is named "memory" and the second mode $b$ named "buffer": the buffer is leaky and a 2-to-1 photon conversion coupling between the memory and the buffer is implemented \cref{fig_diagram_cat_architecture}. For the photon conversion, a three-wave mixing process, to work we need to account for energy conservation: we opted for the autoparametric architecture \cite{Marquet2024Apr} where the frequency of the buffer is tailored to match twice the frequency of the memory, but one is free to add some pump to account for any frequency disparities \TMD{Add citations on pumped cat from Alice \& Bob and Quantic ENS team}. The reason for this choice is that the autoparametric cat is a Floquet system and is thus easier to analyze.

\begin{figure}[tbp]
    \centering
    \begin{subfigure}[t]{0.34\columnwidth}
        \centering
        \includegraphics[width=\columnwidth]{Figures/cat-assets/Feynman_diagram_one_mode.pdf}
        \caption{Diagram of the one-mode cat}
    \end{subfigure}
    \hfill
    \begin{subfigure}[t]{0.64\columnwidth}
        \centering
        \includegraphics[width=\columnwidth]{Figures/cat-assets/Feynman_diagram_two_mode.pdf}
        \caption{Diagram of the autoparametric two-mode cat}
    \end{subfigure}
    \caption{Diagrams of different cat architectures}
    \label{fig_diagram_cat_architecture}
\end{figure}

% Requires in the preamble: \usepackage{circuitikz} and the colors bufcol, cqcol, gcol.
% \resizebox scales the whole drawing (lines, components, text) to the column width.
\begin{figure*}[tbp]
  \centering
\resizebox{\columnwidth}{!}{%
\begin{circuitikz}[line cap=round,line join=round,
  % thickness=1: coil drawn with exactly the same width as the wire it sits on
  circuit_L/.style = {cute inductor, inductors/scale=1.5,
    inductors/width=.5, inductors/coils=4, inductors/thickness=1},
  external_coil/.style = {cute inductor, inductors/scale=1.5,
    inductors/width=.7, inductors/coils=4, inductors/thickness=1}]
  % ------------------------------------------------
  % Global styles (fonts are scaled down by \resizebox:
  % 16pt -> ~9pt, 18pt -> ~10pt in the final column)
  % ------------------------------------------------
  \tikzset{
    buffer/.style={draw=bufcol, text=bufcol, line width=1.6pt},
    catq/.style={draw=cqcol, text=cqcol, line width=1.6pt},
    gstyle/.style={draw=gcol, text=gcol, line width=1.6pt},
    labbuf/.style={text=bufcol, font=\fontsize{16}{16}\selectfont},
    titlebuf/.style={text=bufcol, font=\fontsize{18}{18}\selectfont},
    titlecq/.style={text=cqcol, font=\fontsize{18}{18}\selectfont},
    jj/.style={draw=bufcol, line width=2.6pt},
    ccoup/.style={line width=2.6pt}
  }

  % ------------------------------------------------
  % Coordinates
  % ------------------------------------------------
  \coordinate (BLB) at (0,0);
  \coordinate (BLT) at (0,4.2);
  \coordinate (BTopL) at (2.6,4.2);
  \coordinate (BTopR) at (5.2,4.2);
  \coordinate (BRB) at (5.2,0);
  \coordinate (CapLeft)  at (6.85,4.62);
  \coordinate (CapRight) at (7.15,4.62);
  \coordinate (QLT) at (8.7,4.2);
  \coordinate (QLB) at (8.7,0);
  \coordinate (QRT) at (12.0,4.2);

  % ------------------------------------------------
  % Buffer outer frame and top wire
  % ------------------------------------------------
  \draw[buffer] (BLB) -- (BLT) -- (BTopL);
  \draw[buffer] (BTopL) -- (2.6,4.62) -- (CapLeft);
  \draw[buffer] (BTopL) -- (BTopR) -- (BRB) -- (BLB);
  \draw[buffer] (2.6,4.62) -- (0.875,4.62);

  % ------------------------------------------------
  % Memory frame (inductor inherits the 1.6pt wire width)
  % ------------------------------------------------
  \draw[catq] (CapRight) -- (10.35,4.62) -- (10.35,4.2);
  \draw[catq] (QLT) -- (QRT) to[C, -] ++(0,-4.2) -- (QLB)
              to[circuit_L, -] ++(0,4.2) -- cycle;

  % ------------------------------------------------
  % Coupling capacitors
  % ------------------------------------------------
  \draw[buffer, ccoup] (0.875,4.62) -- ++(0,-0.28) -- ++(0,0.56);
  \draw[gstyle, ccoup] (0.575,4.62) -- ++(0,-0.28) -- ++(0,0.56);
  \draw[buffer, ccoup] (CapLeft)  -- ++(0,-0.28) -- ++(0,0.56);
  \draw[catq,   ccoup] (CapRight) -- ++(0,-0.28) -- ++(0,0.56);

  % ------------------------------------------------
  % Buffer: Josephson junctions, inductor and capacitor
  % ------------------------------------------------
  \draw[jj] (-0.22,2.42) -- (0.22,1.88);
  \draw[jj] (-0.22,1.88) -- (0.22,2.42);
  \draw[buffer] (3.45,4.2) to[C, -] ++(0,-4.2);
  \draw[buffer] (1.75,0) to[circuit_L, -] ++(0,4.2);
  \draw[jj] (4.98,2.42) -- (5.42,1.88);
  \draw[jj] (4.98,1.88) -- (5.42,2.42);

  % ------------------------------------------------
  % Grounds
  % ------------------------------------------------
  \draw[buffer] (2.6,0) -- (2.6,-0.35);
  \draw[buffer] (2.30,-0.35) -- (2.60,-0.65) -- (2.90,-0.35) -- cycle;
  \draw[catq] (10.35,0) -- (10.35,-0.35);
  \draw[catq] (10.05,-0.35) -- (10.35,-0.65) -- (10.65,-0.35) -- cycle;

  % ------------------------------------------------
  % External flux lines (one continuous path each)
  % ------------------------------------------------
  \draw[gstyle] (0.1,-2) -- (0.1,-0.5) to[external_coil, -] (1.65,-0.5) -- (1.65,-2);
  \draw[gstyle] (3.6,-2) -- (3.6,-0.5) to[external_coil, -] (5.15,-0.5) -- (5.15,-2);

  % ------------------------------------------------
  % Input line (coax + ground)
  % ------------------------------------------------
  \draw[gstyle] (0.575,4.62) -- (-1,4.62);
  \draw[gstyle] (-1,4) -- (-3,4);
  \draw[gstyle] (-1,5.24) -- (-3,5.24);
  \draw[gstyle] (-2,4) -- (-2,3.65);
  \draw[gstyle] (-2.3,3.65) -- (-2,3.35) -- (-1.7,3.65) -- cycle;
  \draw[gstyle] (-1,4.62) ellipse (0.2 and 0.62);

  % ------------------------------------------------
  % Titles and labels
  % ------------------------------------------------
  \node[titlebuf] at (2.6,5.0)   {Buffer};
  \node[titlecq]  at (10.35,5.0) {Memory};
  \node[labbuf]   at (-0.72,2.15) {$E_{J}$};
  \node[labbuf]   at (5.90,2.15)  {$E_{J}$};
\end{circuitikz}%
}
\caption{Example of circuit implementing the autoparametric cat qubit.}
\label{fig:circuit_basic}
\end{figure*}
Such qubit is built with the three musketeers of superconducting circuits: capacitor, inductors and Josephson junctions. In an ideal world, this last element brings all the non-linearity in the system thanks to the relation between voltage $V$ and phase $\phi$
\begin{equation}
    V \propto \sin(\phi/\Phi_0),
\end{equation}
where $\Phi_0=\hbar/2e$ is the reduced magnetic flux quantum. The Hamiltonian of such a circuit contains transcendental (trigonometric) terms and any polynomial expression is only an approximation.

In this paper, we researchers living in an ideal world, we focus on cat qubit without any parasitic terms, but due to some personal repulsion for approximations, we will keep the whole nonlinearity.

The Hamiltonian of a dissipative cat qubit reads \cite{Carde2025Feb}
\begin{equation}
\begin{split}
    H =\;& \omega_a\ha\da\ha+\omega_b\hb\da\hb+\eps_d\cos(\omega_d t)(\hb+\hb\da)\\
    &-g\sin(\hat\varphi_a+\hat\varphi_b+\varphi_d).
\end{split}
\label{eq_ham_0}
\end{equation}
We opted the flux $\varphi_d$ such that we cancel the linear term of the sine coupling
\begin{equation}
\begin{split}
    H =\;& \omega_a\ha\da\ha+\omega_b\hb\da\hb+\eps_d\cos(\omega_d t)(\hb+\hb\da)\\
    &-g\bigl[\sin(\hat\varphi_a+\hat\varphi_b)-\hat\varphi_a-\hat\varphi_b\bigr],
\end{split}
\label{eq_ham}
\end{equation}
to minimize the frequency shift as the coupling rate $g$ is varied.
The dissipation is also considered as ideal, meaning that there is only single-photon loss on the buffer
\begin{equation}
    L=\kappa_b\mathcal{D}[b].
\end{equation}
Previous works have investigated the impact of loss on the memory on the dynamics\cite{Carde2025Feb} or on the bit-flip time\cite{Dubovitskii2025Jan}.

\section{Classical analysis}
\label{sec:classical_analysis}
In order to build some intuition, we first look at the classical dynamics of \cref{eq_ham}. There are two ways to derive the equations that we studied. The first one consist of considering the quantum Hamiltonian as a classical Hamiltonian by asking the operators to kindly doff their hats, then we derive Hamilton equation and we add decay on one of the quadrature.
The second way is to start with the GKSL-equation, derive the equation on the Husimi Q-function, approximate it by an isotropic normal gaussian and write down the equation on the mean (\cref{app:derivation_classical_equation}).
\begin{equation}
    \begin{cases}
        \dot a &= -i\omega_a a+ig\,\varphi_a\bigl(\cos X-1\bigr),\\
        \dot b &= -\Bigl(i\omega_b+\frac{\kappa_b}{2}\Bigr)b-i\eps_d\cos(\omega_d t)\\
        &\quad+ig\,\varphi_b\bigl(\cos X-1\bigr),
    \end{cases}
\label{eq_classical_dynamics}
\end{equation}
with $X=\varphi_a(a+a^*)+\varphi_b(b+b^*)$.

\subsection{Phase-space trajectories}
\label{subsec:phase_space_traj}
The first analysis consists of launching many trajectories in phase-space and study their behavior: do some diverge? do some converge? The trajectories are shown on \cref{fig_limit_cycles}.Should we consider the RWA cat, all the trajectories ends up either in $\ket{\alpha,0}$ or $\ket{-\alpha,0}$. Here we take all the non-linearities and rotating terms into account, so that the dynamics is more involved. First of all, in addition to the trajectories (in blue) converging close to $\ket{\pm\alpha,0}$, some trajectories (in red) converge towards excited limit cycles named hereafter the satellite. Due to the transcendental (sine) coupling in the Hamiltonian \cref{eq_ham}, there are infinitely many limit cycles that are solution of the dynamical system \cref{eq_classical_dynamics}. The satellite presented on \cref{fig_limit_cycles} are the first excited stable ones. To summarize, we got two sets of limit cycles: the blue ones corresponding to the cat states and the red satellites. The question becomes now: how do these satellites impact the dynamics that generates the blue limit-cycles.

\begin{figure}[tbp]
    \centering
    \includegraphics[width=0.5\columnwidth]{Chaos/figures_article/classical_limit_cycles_HB_12_10_2026.pdf}
    \caption{Limit cycles for the buffer phase space (left) and the memory phase space (right): blue dots corresponds to the cat limit cycles, found by letting trajectories converge to them. The red dots represent the satellite limit cycles, also found by letting trajectories converge in phases space. The yellow dashed line corresponds to the prediction of the harmonic balance.}
    \label{fig_limit_cycles}
\end{figure}

\subsection{Study of stability of the limit-cycles : Hill's method}
\label{subsec:hill}
In order to answer this question, we study the stability of the satellites thanks to Hill's method and follow the steps described in \cite{Detroux2015Nov}. We first apply harmonic balance and find an approximation of the satellite limit cycle (yellow dashed lines on \cref{fig_limit_cycles}): the error on the approximation is controlled by the number of harmonics that are kept in the Ansatz for the harmonic balance, the larger, the more precise. Note that you can also derive analytically the center and the radii of the satellite \cref{app:derivation_harmonic_balance}.

To study the stability of the satellite limit-cycle, we focus on the dynamics at the vicinity of them, like we would do for a fixed point. 
Let $\boldsymbol{x}^*(t)$ be the periodic solution of a limit-cycle, we perturb this solution
\begin{equation}
    \boldsymbol{x}(t)=\boldsymbol{x}^*(t)+e^{\mu t}\boldsymbol{\delta}\boldsymbol{x}(t).
\end{equation}

Since we work in a 4d space, they are 4 complex numbers describing the dynamics of the perturbation $\left\{e^{\mu_iT}\right\}_{i\in\{1,2,3,4\}}$: they are called the Floquet multipliers\footnote{Note that the Floquet multipliers consists either of two pairs of complex conjugate numbers or of 4 real coefficients, so that the dynamics is effectively described by $n$ real numbers only in $n$-d.} There are three scenarii to take into considerations
\begin{itemize}
    \item $\forall i \: |e^{\mu_iT}|<1$: nearby trajectories spiral towards the limit cycle, it is stable
    \item $\exists i\:|e^{\mu_iT}|>1$:nearby trajectories spiral away of the limit cycle, it is unstable
    \item $\exists i\:e^{\mu_iT}=-1$:the limit cycle experience period doubling, it becomes unstable but a new limit-cycle with a period twice as big appears and its stability should be investigated
\end{itemize}

The first parameter we tune is the two-photon exchange rate $g_2$. For low value of $g_2$, the limit cycle is stable, then it undergoes period-doubling as $g_2$ is increased and becomes unstable for large value. A new cycle has appeared due to the period-doubling, which is stable at first but quickly becomes unstable as well. The period-2 orbit actually loses stability via a Neimark–Sacker bifurcation (complex pair crossing $|e^{\mu_i T}|=1$). \TMD{Link NS bifurcation and chaos in previous works?} We now show that this instability threshold is related to the onset of chaos in this system.


\begin{figure}[tbp]
    \centering
    % --- Left column: three small plots stacked vertically ---
    \begin{minipage}[c]{0.2\columnwidth}
        \centering
        \begin{subfigure}[tbp]{\columnwidth}
            \centering
            \includegraphics[width=\columnwidth]{Figures/article_chaos/stabilization_limit_cycle.pdf}
            \caption{Stabilization}
            \label{fig_scheme_stabilization}
        \end{subfigure}
        \vspace{0.5cm}
        \begin{subfigure}[tbp]{\columnwidth}
            \centering
            \includegraphics[width=\columnwidth]{Figures/article_chaos/period_doubling_limit_cycles.pdf}
            \caption{Period doubling}
            \label{fig_scheme_period_doubling}
        \end{subfigure}
        \vspace{0.5cm}
        \begin{subfigure}[tbp]{\columnwidth}
            \centering
            \includegraphics[width=\columnwidth]{Figures/article_chaos/instability_limit_cycles.pdf}
            \caption{Instability}
            \label{fig_scheme_unstabilization}
        \end{subfigure}
    \end{minipage}
    \hfill
    % --- Right column: two plots stacked vertically ---
    \begin{minipage}[c]{0.78\columnwidth}
        \centering
        \begin{subfigure}[tbp]{\columnwidth}
            \centering
            \includegraphics[width=\columnwidth]{Chaos/figures_article/floquet_multiplier_modulus_before_period_doubling_12_10_2026.pdf}
            \caption{Below period doubling threshold: $T=2\pi/\omega_b$}
            \label{fig_floq_mult_below_period_doubling}
        \end{subfigure}
        \vspace{0.5cm}
        \begin{subfigure}[tbp]{\columnwidth}
            \centering
            \includegraphics[width=\columnwidth]{Chaos/figures_article/floquet_multiplier_modulus_after_period_doubling_12_10_2026.pdf}
            \caption{Above period doubling threshold: $T=4\pi/\omega_b$}
            \label{fig_floq_mult_above_period_doubling}
        \end{subfigure}
    \end{minipage}
    \caption{Study of the stability of the satellite limit-cycles.}
    \label{fig_hill_method}
\end{figure}


\subsection{Characterizing chaos}
\label{subsec:chaos_lyap}

The most accessible signature of chaos is to analyze phase space trajectories as done on \cref{fig_classical_onset_of_chaos_trajectories_low,fig_classical_onset_of_chaos_trajectories_mid,fig_classical_onset_of_chaos_trajectories_high}. For low value of $g_2$, all trajectories in phase space are integrable (cf \cref{fig_classical_onset_of_chaos_trajectories_low}), then when $g_2$ exceeds the stability threshold, two chaotic basins develop around the satellites (cf \cref{fig_classical_onset_of_chaos_trajectories_mid}) and finally, if $g_2$ keeps increasing, these basins grow and take over the phase space (cf \cref{fig_classical_onset_of_chaos_trajectories_high}).

In order to properly characterize the classical chaos, we use the maximal Lyapunov exponent. We run into the issue that a Lyapunov exponent is defined for a given trajectory in phase space, \textit{i.e.} a given initial point. We chose to focus on the trajectory initialized in the phase-space origin $(0,0)$. The reason is that a cat state is roughly composed of two features: two lobes and fringes in between them. The two lobes corresponds to attractors in phase space, so we expect the fringes to be more impacted by the chaotic basins. The fringes are centered on the origin, so we chose to focus on trajectories starting on this point.

On \cref{fig_classical_onset_of_chaos_heatmap}, we tune two parameters: the two-photon exchange rate $g_2$ on the x axis and the buffer drive $\eps_d$, which determine the cat size (when a cat can be defined). This heatmap can be separated in three blocs:
\begin{itemize}
    \item The blue region on the left (with point a), where the trajectory initialized at the origin is regular
    \item The blue region in the middle (with point b, the chaotic basin develop around the satellite limit-cyles but the trajectory that starts in the origin stay regular
    \item The (vastly) red region on the right (with point c), the trajectory is (almost) always chaotic \footnote{The blue islands are due to highly excited stable limit-cycle: due to chaos, the trajectory starting in the origin is quickly sent to these stable cycles and it appears as regular}.
\end{itemize}

We have found an upper threshold in some classical limit of the autoparametric dissipative cat, can we do the same on the quantum model directly?

\begin{figure}[tbp]
    \centering
    % --- Left column: three small plots stacked vertically ---
    \begin{minipage}[c]{0.2\columnwidth}
        \centering
        \begin{subfigure}[tbp]{\columnwidth}
            \centering
            \includegraphics[width=\columnwidth]{Chaos/figures_article/phase_space_traj_g2_60MHz_12_10_2026.png}
            \caption{Low $g_2$.}%: regular, periodic trajectory.}
            \label{fig_classical_onset_of_chaos_trajectories_low}
        \end{subfigure}
        \vspace{0.5cm}
        \begin{subfigure}[tbp]{\columnwidth}
            \centering
            \includegraphics[width=\columnwidth]{Chaos/figures_article/phase_space_traj_g2_90MHz_12_10_2026.png}
            \caption{Intermediate $g_2$.}%: chaotic basins start to develop.}
            \label{fig_classical_onset_of_chaos_trajectories_mid}
        \end{subfigure}
        \vspace{0.5cm}
        \begin{subfigure}[tbp]{\columnwidth}
            \centering
            \includegraphics[width=\columnwidth]{Chaos/figures_article/phase_space_traj_g2_105MHz_12_10_2026.png}
            \caption{High $g_2$.}%: fully chaotic trajectories.}
            \label{fig_classical_onset_of_chaos_trajectories_high}
        \end{subfigure}
    \end{minipage}
    \hfill
    % --- Right column: the single large plot ---
    \begin{subfigure}[c]{0.78\columnwidth}
        \centering
        \includegraphics[width=\columnwidth]{Chaos/figures_article/lyapunov_heatmap_origin_start_12_10_2026.pdf}
        \caption{Heatmap of the maximum Lyapunov exponent for the trajectory starting in the phase space origin $(a=0, b=0)$. The yellow line is just there to separate the regular region (blue) from the chaotic one (red)}
        \label{fig_classical_onset_of_chaos_heatmap}
    \end{subfigure}
    \caption{Classical onset of chaos}
    \label{fig_classical_onset_of_chaos}
\end{figure}

\section{Quantum analysis}
\label{sec:quantum_analysis}

A difficulty that arises in the quantum model is to find the right metric for quantum chaos. Quantum chaos in open quantum systems has been investigated under diverse lenses: inverse participation ratio \cite{Borner2024Aug}, observation entropy \cite{Pg2022Dec}, truncated Wigner method \cite{Pappalardi2020Aug}, spectral properties of the Lindbladian \cite{Sa2020Apr}. In order to be able to compare the quantum results with the classical prediction, we chose to work with a metric that provides a quantum mechanical analog to the maximal Lyapunov exponent: the Out-Of-Time-Order Correlator (OTOC) \cite{Larkin1969Jun}. The OTOC is defined as 
\begin{equation}
    C(t)=-\left\langle\left[\hat x_a(t),\hat p_a(0)\right]^2\right\rangle
\end{equation}
and at short time, $C(t)\sim e^{2\lambda t}$, where $\lambda$ is the maximal Lyapunov exponent when we take the quasiclassical limit of the commutator \cite{Garcia-Mata2022Sep}.

We propagate the position operator over time with the adjoint Lindblad master equation \cite[p.~125]{Breuer2006Jan}
\begin{equation}
    \dot{\hat x}_a(t) = i[H(t),\hat x_a(t)]+\kappa_b \left(\hat b\da\hat x_a(t)\hat b -\frac 1 2 \{\hat b\da \hat b, \hat x_a(t)\}\right),
\end{equation},
compute the OTOC over an approximation of the stabilized cat-state and fit the OTOC as an exponential.

 First, we see two regimes: at low values of the two-photon exchange rate $g_2$, the cat is stabilized (cf \cref{fig_wigner_cat_low}) and for large values of $g_2$, the Wigner distribution is delocalized over the phase space (cf \cref{fig_wigner_cat_high}). The evolution of the OTOC is showed in the inset of \cref{fig_quantum_onset_of_chaos}: for small $g_2$, the OTOC stays flat and for large $g_2$, it grows like an exponential with some modulation (since the Lindbladian is time-dependent\footnote{The curve is not periodic because the system is chaotic.}). 

Finally, we draw heatmap of the Lyapunov exponent extracted from the OTOC where we again act on the two-photon exchange rate $g_2$ and the buffer drive $\eps_d$. We again find that increasing either parameter drive the system from a regular dynamics to chaos. Nevertheless, the threshold in the quantum model is below the classical threshold: Lyapunov exponent from the OTOC are positive even below the yellow line on \cref{fig_quantum_onset_of_chaos}. This is explained by the intrinsic quantum delocalization: classical dynamics can be oblivious to the chaotic basin whereas the quantum dynamics can be more easily impacted by them.

One questions remains though: are there ways to push the threshold to larger parameters range?

\begin{figure}[h]
    \centering
    \begin{minipage}[c]{0.25\columnwidth}
    \centering
    \begin{subfigure}[tbp]{\columnwidth}
        \centering
        \includegraphics[width=\columnwidth]{Figures/Classical_chaos/CSI_report/wigner_cat_stable.pdf}
        \caption{Low $g_2$.}%: regular, periodic trajectory.}
        \label{fig_wigner_cat_low}
    \end{subfigure}
    \vspace{0.5cm}
    \begin{subfigure}[tbp]{\columnwidth}
        \centering
        \includegraphics[width=\columnwidth]{Figures/Classical_chaos/CSI_report/wigner_cat_unstable.pdf}
        \caption{High $g_2$.}%: chaotic basins start to develop.}
        \label{fig_wigner_cat_high}
    \end{subfigure}
    \end{minipage}
    \hfill
    \begin{subfigure}[c]{0.7\columnwidth}
        \centering
        \includegraphics[width=\columnwidth]{Chaos/figures_article/Lyapunov_from_OTOC_init_cat_12_10_2026.pdf}
        \caption{Heatmap of the maximum Lyapunov exponent found by fitting the OTOC. The average is taken on a cat-state.}
        \label{fig_quantum_onset_of_chaos_heatmap}
    \end{subfigure}
    \caption{Quantum onset of chaos. The subfigures \cref{fig_wigner_cat_low} and \cref{fig_wigner_cat_high} represent Wigner distribution obtained after some times when the initial state is a cat state.}
    \label{fig_quantum_onset_of_chaos}
\end{figure}


\section{Mitigation techniques}
\label{sec:mitigation}
In this last section, we present three ways to keep the regular dynamics on larger range.

\subsection{Increasing the single photon loss}
\label{subsec:increasing_photon_loss}

A first solution, which allows one to keep the initial circuit, would be to increase the single photon loss rate $\kappa_b$ on the buffer. The instability threshold increases linearly with the single photon rate \TMD{Explanation?}, at large value of the rate, cf \cref{fig:instability_vary_kappa_b}.

\begin{figure}[tbp]
    \centering
    \includegraphics[width=\columnwidth]{Chaos/figures_article/instability_vary_kappa_b_12_10_2026.pdf}
    \caption{Study of the instability threshold when the single photon rate on the buffer is varied. The gray dashed line corresponds to the value taken for the study in hereinabove plots.}
    \label{fig:instability_vary_kappa_b}
\end{figure}

\subsection{Multi-buffer architecture}
\label{subsec:multi_buffer}
In a similar spirit, instead of increasing the single photon loss rate, we could maybe introduce new loss channels to mitigate the apparition of chaos?
The second solution is to increase the number of buffers, so that the energy is split between them. The Hamiltonian of such a system reads
\begin{equation}
    \begin{split}
        H_\mathrm{m.b.} =\;& \om_a\hat a^\dagger \hat a + \sum\biggl\{\om_{b_n} \hat b_n^\dagger \hat b_n + \eps_{d_n}\cos(\om_{d_n}t)(\hat b_n + \hat b_n\da) \\&-g_n\left[\sin\left(\hat \varphi_a + \hat  \varphi_{b_n} \right)-\hat \varphi_a - \hat \varphi_{b_n}\right]\biggr\}.
    \end{split}
    \label{eq:multi-buffer-ham}
\end{equation}
Shall we perform RWA on this system, we would obtain
\begin{equation}
    H_\mathrm{m.b.}^\mathrm{RWA}= \sum \eps_{d_n}(\hat b_n + \hat b_n\da)/2 + \left( a^2 \sum_n g_{n,2}  b_n^\dagger + \text{h.c.} \right),
\end{equation}
hence we can define some effective two-photon exchange rate 
\begin{equation}
    g_2^\mathrm{eff}=\sum_n g_{n,2}.
\end{equation}
We compute the maximal Lyapunov for one, two and three buffers which supposedly have the same frequency $\om_{b_n}=2\om_a$ and identical drive amplitude. As shown on \cref{fig:Lyap_multi_buffer_each}, the more buffer we add, the higher the threshold for chaos.
\begin{figure}[tbp]
    \centering
    \includegraphics[width=\columnwidth]{Chaos/figures_article/lyapunov_multi_buffer_12_10_2026.pdf}
    \caption{Maximum Lyapunov exponent $\lambda$ for the trajectory initiated in the phase-space origin}
    \label{fig:Lyap_multi_buffer_each}
\end{figure}

An example of a circuit that would implement the multi-buffer Hamiltonian \cref{eq:multi-buffer-ham} is presented on \cref{fig_multibuffer_circuit}. However such circuit can be quite painful to design and work with since you increase the number of modes to fabricate and to calibrate afterwards. Couldn't we find another solution?

\begin{figure*}[tbp]
  \centering
\resizebox{\textwidth}{!}{%
\begin{circuitikz}[line cap=round,line join=round,
  circuit_L/.style = {cute inductor, inductors/scale=1.5,
    inductors/width=.5, inductors/coils=4, inductors/thickness=1},
  buffer/.style={draw=bufcol, text=bufcol, line width=1.6pt},
  catq/.style={draw=cqcol, text=cqcol, line width=1.6pt},
  gstyle/.style={draw=gcol, text=gcol, line width=1.6pt},
  labbuf/.style={text=bufcol, font=\fontsize{16}{16}\selectfont},
  titlebuf/.style={text=bufcol, font=\fontsize{18}{18}\selectfont},
  titlecq/.style={text=cqcol, font=\fontsize{18}{18}\selectfont},
  jjbuf/.style={draw=bufcol, line width=2.6pt},
  ccoup/.style={line width=2.6pt}]

  % ================================================
  % MEMORY (LC resonator), x in [-6.4, -3.1]
  % ================================================
  \draw[catq] (-6.4,0) to[circuit_L, -] (-6.4,4.2) -- (-3.1,4.2);
  \draw[catq] (-3.1,4.2) to[C, -] (-3.1,0) -- (-6.4,0);
  \draw[catq] (-4.75,0) -- (-4.75,-0.35);
  \draw[catq] (-5.05,-0.35) -- (-4.75,-0.65) -- (-4.45,-0.35) -- cycle;
  \node[titlecq] at (-4.75,-1.2) {Memory $a$};

  % ================================================
  % MEMORY BUS (blue), dashed between buffer 2 and buffer N
  % ================================================
  \draw[catq] (-4.75,4.2) -- (-4.75,6.2) -- (14.4,6.2);
  \draw[catq, dashed] (14.4,6.2) -- (18.8,6.2);
  \draw[catq] (18.8,6.2) -- (22.6,6.2);

  % ================================================
  % BUFFERS: ATS (JJ, L, C, JJ) + coupling to bus + own grey port
  % \x = left edge of the ATS frame, \n = buffer index
  % ================================================
  \foreach \x/\n in {0/1, 9/2, 20/N}{
    % ATS frame: junction | inductor | capacitor | junction
    \draw[buffer] (\x,0) -- (\x,4.2) -- (\x+5.2,4.2) -- (\x+5.2,0) -- cycle;
    \draw[buffer] (\x+1.75,0) to[circuit_L, -] ++(0,4.2);
    \draw[buffer] (\x+3.45,4.2) to[C, -] ++(0,-4.2);
    \draw[jjbuf] (\x-0.22,2.42) -- (\x+0.22,1.88);
    \draw[jjbuf] (\x-0.22,1.88) -- (\x+0.22,2.42);
    \draw[jjbuf] (\x+4.98,2.42) -- (\x+5.42,1.88);
    \draw[jjbuf] (\x+4.98,1.88) -- (\x+5.42,2.42);
    % ground
    \draw[buffer] (\x+2.6,0) -- (\x+2.6,-0.35);
    \draw[buffer] (\x+2.3,-0.35) -- (\x+2.6,-0.65) -- (\x+2.9,-0.35) -- cycle;
    \node[titlebuf] at (\x+2.6,-1.2) {Buffer $b_{\n}$};
    % node above the ATS, coupling capacitor to the memory bus
    \draw[buffer] (\x+2.6,4.2) -- (\x+2.6,5.27);
    \draw[buffer, ccoup] (\x+2.6,5.27) -- ++(-0.28,0) -- ++(0.56,0);
    \draw[catq,   ccoup] (\x+2.6,5.57) -- ++(-0.28,0) -- ++(0.56,0);
    \draw[catq] (\x+2.6,5.57) -- (\x+2.6,6.2);
    \fill[cqcol] (\x+2.6,6.2) circle (0.09);
    % own grey port, capacitively coupled to the buffer node
    \draw[buffer] (\x+2.6,4.62) -- (\x+4.6,4.62);
    \fill[bufcol] (\x+2.6,4.62) circle (0.09);
    \draw[buffer, ccoup] (\x+4.6,4.62) -- ++(0,-0.28) -- ++(0,0.56);
    \draw[gstyle, ccoup] (\x+4.9,4.62) -- ++(0,-0.28) -- ++(0,0.56);
    \draw[gstyle] (\x+4.9,4.62) -- (\x+6.3,4.62);
    \draw[gstyle] (\x+6.3,4.62) ellipse (0.2 and 0.62);
    \draw[gstyle] (\x+6.3,4.0) -- (\x+7.7,4.0);
    \draw[gstyle] (\x+6.3,5.24) -- (\x+7.7,5.24);
    \draw[gstyle] (\x+7.0,4.0) -- (\x+7.0,3.65);
    \draw[gstyle] (\x+6.7,3.65) -- (\x+7.0,3.35) -- (\x+7.3,3.65) -- cycle;
  }
  % dots standing for buffers 3 ... N-1
  \node[titlebuf] at (16.6,2.1) {$\cdots$};
\end{circuitikz}%
}
  \caption{Example of a multi-buffer circuit implementing $H_\mathrm{m.b.}$.}
  \label{fig_multibuffer_circuit}
\end{figure*}

\subsection{Parity-protected cat qubit}
\label{subsec:parity_protected}
Should the recommendation from this article be condensed in one sentence, this would likely be: if you seek for less chaos, remove as many harmonics in your Hamiltonian as you can. This is what we demonstrate here: we focus on the design presented on \TMD{Cite Antoine Marquet?} whose Hamiltonian reads
\begin{equation}
\begin{split}
    H_\mathrm{p.p.} =\;&  \om_aa^\dagger a + \om_{b} b^\dagger b+\eps_d\cos(\omega_d t)(\hb+\hb\da)\\& -g\left[\cos\left(\hat \varphi_a\right)\sin\left(\hat\varphi_{b}\right)-\hat\varphi_{b}\right].
\label{eq_ham_pp}
\end{split}
\end{equation}
We call this the parity-protected version of our initial Hamiltonian \cref{eq_ham}: the sine coupling is replaced with this product of cosine and sine, such that, under the change $\hat a \to -\hat a$, $H_\mathrm{p.p.}$ stays the same, just as with $H^\mathrm{RWA}$.

This Hamiltonian has also satellite, as shown on \cref{fig:phase_space_traj_parity_protected}, but they don't undergo any period doubling. Indeed, none of the Floquet multiplier never reaches $-1$ as shown on \cref{fig:floq_mult_parity_protected}.

\begin{figure}[tbp]
    \centering
    % Row 1
    \begin{subfigure}{0.45\columnwidth}
        \centering
        \includegraphics[width=\linewidth]{Chaos/figures_article/parity_protected_phase_space_traj_g2_100MHz_12_10_2026.png}
        \caption{Phase-space limit cycles}
        \label{fig:phase_space_traj_parity_protected}
    \end{subfigure}
    \hfill
    \begin{subfigure}{0.45\columnwidth}
        \centering
        \includegraphics[width=\linewidth]{Chaos/figures_article/floquet_multiplier_parity_protected_12_10_2026.pdf}
        \caption{Real part of the Floquet multipliers}
        \label{fig:floq_mult_parity_protected}
    \end{subfigure}

    \caption{Analysis of the parity protected cat qubit}
\label{fig:parity_protected}
\end{figure}

A circuit that implements such parity-protected Hamiltonian is drawn on \cref{fig:app_circuit}: its mirrored architecture allows to keep parity symmetry safe, the derivation of the Hamiltonian is presented in \cref{app:parity_circuit}. 

% Requires in the preamble: \usepackage{circuitikz} and the colors bufcol, cqcol, gcol.
% \resizebox scales the whole drawing (lines, components, text) to the column width.
\begin{figure}[tbp]
  \centering
\resizebox{\columnwidth}{!}{%
\begin{circuitikz}[line cap=round,line join=round,
  circuit_L/.style = {cute inductor, inductors/scale=1.5,
    inductors/width=.5, inductors/coils=4, inductors/thickness=1},
  external_coil/.style = {cute inductor, inductors/scale=1.5,
    inductors/width=.7, inductors/coils=4, inductors/thickness=1}]
  % ------------------------------------------------
  % Global styles (fonts are scaled by \resizebox, ~0.45:
  % 20pt -> ~9pt, 22pt -> ~10pt in the final column)
  % ------------------------------------------------
  \tikzset{
    buffer/.style={draw=bufcol, text=bufcol, line width=1.6pt},
    catq/.style={draw=cqcol, text=cqcol, line width=1.6pt},
    gstyle/.style={draw=gcol, text=gcol, line width=1.6pt},
    labbuf/.style={text=bufcol, font=\fontsize{20}{20}\selectfont},
    labcq/.style={text=cqcol, font=\fontsize{20}{20}\selectfont},
    titlebuf/.style={text=bufcol, font=\fontsize{22}{22}\selectfont},
    titlecq/.style={text=cqcol, font=\fontsize{22}{22}\selectfont},
    jjbuf/.style={draw=bufcol, line width=2.6pt},
    jjcq/.style={draw=cqcol, line width=2.6pt},
    ccoup/.style={line width=2.6pt}
  }

  % ================================================
  % BUFFER (ATS + capacitor), left block: x in [0, 5.2]
  % ================================================
  \draw[buffer] (0,0) -- (0,4.2) -- (5.2,4.2) -- (5.2,0) -- cycle;
  \draw[buffer] (2.6,4.2) -- (2.6,4.62);
  \draw[buffer] (0.875,4.62) -- (6.85,4.62);            % top wire
  \draw[buffer] (3.45,4.2) to[C, -] ++(0,-4.2);          % capacitor
  \draw[buffer] (1.75,0) to[circuit_L, -] ++(0,4.2);     % inductor
  \draw[jjbuf] (-0.22,2.42) -- (0.22,1.88);  \draw[jjbuf] (-0.22,1.88) -- (0.22,2.42);
  \draw[jjbuf] (4.98,2.42) -- (5.42,1.88);   \draw[jjbuf] (4.98,1.88) -- (5.42,2.42);
  \draw[buffer] (2.6,0) -- (2.6,-0.35);
  \draw[buffer] (2.30,-0.35) -- (2.60,-0.65) -- (2.90,-0.35) -- cycle;

  % ================================================
  % MEMORY: copy of the ATS + capacitor, x in [8.7, 13.9]
  % ================================================
  \draw[catq] (8.7,0) -- (8.7,4.2) -- (13.9,4.2) -- (13.9,0) -- cycle;
  \draw[catq] (11.3,4.2) -- (11.3,4.62);
  \draw[catq] (7.15,4.62) -- (13.3,4.62) -- (13.3,4.97); % top wire + stub to grey coupler
  \draw[catq] (12.15,4.2) to[C, -] ++(0,-4.2);           % capacitor
  \draw[catq] (10.45,0) to[circuit_L, -] ++(0,4.2);      % inductor
  \draw[jjcq] (8.48,2.42) -- (8.92,1.88);    \draw[jjcq] (8.48,1.88) -- (8.92,2.42);
  \draw[jjcq] (13.68,2.42) -- (14.12,1.88);  \draw[jjcq] (13.68,1.88) -- (14.12,2.42);
  \draw[catq] (11.3,0) -- (11.3,-0.35);
  \draw[catq] (11.0,-0.35) -- (11.3,-0.65) -- (11.6,-0.35) -- cycle;

  % ================================================
  % Coupling capacitors
  % ================================================
  % grey input line -- buffer
  \draw[gstyle, ccoup] (0.575,4.62) -- ++(0,-0.28) -- ++(0,0.56);
  \draw[buffer, ccoup] (0.875,4.62) -- ++(0,-0.28) -- ++(0,0.56);
  % buffer -- memory
  \draw[buffer, ccoup] (6.85,4.62) -- ++(0,-0.28) -- ++(0,0.56);
  \draw[catq,   ccoup] (7.15,4.62) -- ++(0,-0.28) -- ++(0,0.56);
  % shared grey port -- memory
  \draw[catq,   ccoup] (13.3,4.97) -- ++(-0.28,0) -- ++(0.56,0);
  \draw[gstyle, ccoup] (13.3,5.27) -- ++(-0.28,0) -- ++(0.56,0);

  % ================================================
  % Grey lines
  % ================================================
  % left input line (coax + ground)
  \draw[gstyle] (0.575,4.62) -- (-1,4.62);
  \draw[gstyle] (-1,4) -- (-2.4,4);
  \draw[gstyle] (-1,5.24) -- (-2.4,5.24);
  \draw[gstyle] (-1.7,4) -- (-1.7,3.65);
  \draw[gstyle] (-2.0,3.65) -- (-1.7,3.35) -- (-1.4,3.65) -- cycle;
  \draw[gstyle] (-1,4.62) ellipse (0.2 and 0.62);
  % branch of the same port, running above the buffer to the memory coupler
  \draw[gstyle] (-0.4,4.62) -- (-0.4,6.0) -- (13.3,6.0) -- (13.3,5.27);
  \fill[gcol] (-0.4,4.62) circle (0.09);                % T-junction dot
  % flux lines below the buffer
  %M\draw[gstyle] (0.1,-2) -- (0.1,-0.5) to[external_coil, -] %(1.65,-0.5) -- (1.65,-2);
  %\draw[gstyle] (3.6,-2) -- (3.6,-0.5) to[external_coil, -] %(5.15,-0.5) -- (5.15,-2);

  % ================================================
  % Titles and labels
  % ================================================
  \node[titlebuf] at (2.6,5.05)  {Buffer};
  \node[titlecq]  at (11.3,5.05) {Memory};
  \node[labbuf]   at (-0.75,2.15) {$E_{J}$};
  \node[labbuf]   at (5.95,2.15)  {$E_{J}$};
  \node[labcq]    at (7.95,2.15)  {$E_{J}$};
  \node[labcq]    at (14.65,2.15) {$E_{J}$};
\end{circuitikz}%
}
  \caption{Parity-protected cat qubit.}
  \label{fig:app_circuit}
\end{figure}

\section{Conclusion}
In this paper we showcase a pipeline for the study of chaos in Floquet quantum systems: first go to the classical limits, compute the instability threshold with Hill method, characterize the onset of chaos with the maximal Lyapunov exponent and measure the OTOC in the quantum model. We applied this workflow on a cat qubit, namely the autoparametric dissipative cat qubit and show how such a system can be driven from a regular dynamics into a chaotic one. Finally we present three ways to mitigate this chaotic behavior: increasing the single photon loss rate, adding more loss channels, rethinking the circuit architecture to restore parity symmetry.

\begin{acknowledgments}
The authors thanks the ANRT for the fundings, the members of Quantic team (Inria) and of the simulation team (Alice \& Bob).

\end{acknowledgments}

\appendix
\section{Derivation of the classical equations of motion}
\label{app:derivation_classical_equation}

In this Appendix, we derive \cref{eq_classical_dynamics} from the GKSL equation
\begin{equation}
    \dot{\hat\rho} = -i[H,\hat\rho] + \kappa_b\mathcal{D}[\hb]\hat\rho .
    \label{eq_app_GKSL}
\end{equation}
We write $\hat x\equiv\hat\varphi_a+\hat\varphi_b$, with $\hat\varphi_a=\varphi_a(\ha+\ha\da)$ and $\hat\varphi_b=\varphi_b(\hb+\hb\da)$. The Hamiltonian splits into a linear and a nonlinear part, $H=H_l+H_{nl}$, with
\begin{align}
    H_l &= \omega_a\ha\da\ha+\omega_b\hb\da\hb+g\,\hat x \nonumber\\
        &\quad +\eps_d\cos(\omega_d t)(\hb+\hb\da),\\
    H_{nl} &= -g\sin\hat x .
\end{align}

\subsection{Husimi function and correspondence rules}
We work with the two-mode Husimi function
\begin{equation}
    Q(\alpha,\beta)=\frac{1}{\pi^2}\bra{\alpha,\beta}\hat\rho\ket{\alpha,\beta},
\end{equation}
treating $\alpha,\alpha^*,\beta,\beta^*$ as independent variables. Its moments are anti-normally ordered expectation values, $\int d^2\alpha\,d^2\beta\;\alpha^m(\alpha^*)^n Q=\langle\ha^m(\ha\da)^n\rangle$. Products of $\hat\rho$ with mode operators translate into differential operators acting on $Q$ \TMD{Add citation}:
\begin{equation}
\begin{aligned}
    \ha\hat\rho &\mapsto (\alpha+\partial_{\alpha^*})Q, &\quad \hat\rho\ha\da &\mapsto (\alpha^*+\partial_\alpha)Q,\\
    \ha\da\hat\rho &\mapsto \alpha^*Q, &\quad \hat\rho\ha &\mapsto \alpha Q,
\end{aligned}
\label{eq_app_rules}
\end{equation}
and similarly for $\hb$ with $\beta$. Successive left multiplications compose in order, $\hat O_1\hat O_2\hat\rho\mapsto O_1^{L}O_2^{L}Q$, and right multiplications in reverse order. Below, $\partial_\alpha(AQ)+\text{c.c.}$ is shorthand for $\partial_\alpha(AQ)+\partial_{\alpha^*}(A^*Q)$.

\subsection{Linear part}
For a single mode with $H=\omega\ha\da\ha+F(t)(\ha+\ha\da)$, $F$ real, \cref{eq_app_rules} gives
\begin{align}
    -i[\omega\ha\da\ha,\hat\rho] &\mapsto -i\omega\bigl[\alpha^*\partial_{\alpha^*}-\alpha\partial_\alpha\bigr]Q \nonumber\\
    &\;= \partial_\alpha(i\omega\alpha Q)+\text{c.c.},\\
    -iF[\ha+\ha\da,\hat\rho] &\mapsto -iF\bigl(\partial_{\alpha^*}-\partial_\alpha\bigr)Q \nonumber\\
    &\;= \partial_\alpha(iFQ)+\text{c.c.}
\end{align}
For the dissipator, the three terms of $\mathcal{D}[\hb]\hat\rho$ map onto $(\beta+\partial_{\beta^*})(\beta^*+\partial_\beta)Q$, $-\tfrac12\beta^*(\beta+\partial_{\beta^*})Q$ and $-\tfrac12\beta(\beta^*+\partial_\beta)Q$, which sum to
\begin{equation}
    \kappa_b\mathcal{D}[\hb]\hat\rho \mapsto \partial_\beta\Bigl(\frac{\kappa_b}{2}\beta Q\Bigr)+\text{c.c.}+\kappa_b\partial_\beta\partial_{\beta^*}Q .
\end{equation}
Collecting all terms, the linear part of the dynamics reads
\begin{equation}
    \dot Q\big|_l=\partial_\alpha(A^\alpha Q)+\partial_\beta(A^\beta Q)+\text{c.c.}+\kappa_b\partial_\beta\partial_{\beta^*}Q,
    \label{eq_app_Q_linear}
\end{equation}
with the drift coefficients
\begin{align}
    A^\alpha &= i\omega_a\alpha+ig\varphi_a,\\
    A^\beta &= \Bigl(i\omega_b+\frac{\kappa_b}{2}\Bigr)\beta+i\eps_d\cos(\omega_d t)+ig\varphi_b .
\end{align}

\subsection{Nonlinear part}
Since left (right) multiplications compose in (reverse) order, the rules \eqref{eq_app_rules} extend term by term to any power series $f$ of $\hat x$:
\begin{equation}
    f(\hat x)\hat\rho\mapsto f(D_L)Q,\qquad \hat\rho f(\hat x)\mapsto f(D_R)Q,
\end{equation}
with the first-order differential operators
\begin{align}
    D_L &= \varphi_a(\alpha+\alpha^*+\partial_{\alpha^*})+\varphi_b(\beta+\beta^*+\partial_{\beta^*}),\\
    D_R &= \varphi_a(\alpha+\alpha^*+\partial_{\alpha})+\varphi_b(\beta+\beta^*+\partial_{\beta}).
\end{align}
Hence $-i[H_{nl},\hat\rho]\mapsto ig\,[\sin D_L-\sin D_R]\,Q$, and the full equation for the Husimi function is
\begin{equation}
\begin{split}
    \dot Q =\;& \partial_\alpha(A^\alpha Q)+\partial_\beta(A^\beta Q)+\text{c.c.}\\
    &+\kappa_b\partial_\beta\partial_{\beta^*}Q
    +ig\,[\sin D_L-\sin D_R]\,Q .
\end{split}
\label{eq_app_Q_full}
\end{equation}

Although not needed below, the action of $D_{L}$ can be made explicit. Writing $\partial_{\alpha^*}+\alpha+\alpha^*=e^{-h}\partial_{\alpha^*}e^{h}$ with $h=\alpha\alpha^*+(\alpha^*)^2/2$, and similarly for $\beta$, the exponential $e^{iD_L}$ reduces to a translation of $\alpha^*$ and $\beta^*$:
\begin{equation}
\begin{split}
    e^{iD_L}Q(\alpha,\alpha^*,\beta,\beta^*) =\;& e^{iX-(\varphi_a^2+\varphi_b^2)/2}\\
    &\times Q(\alpha,\alpha^*+i\varphi_a,\beta,\beta^*+i\varphi_b),
\end{split}
\end{equation}
with $X=\varphi_a(\alpha+\alpha^*)+\varphi_b(\beta+\beta^*)$, and $e^{iD_R}Q$ is obtained by exchanging the roles of $\alpha$ and $\alpha^*$ (resp.\ $\beta$ and $\beta^*$). The nonlinearity thus couples $Q$ at points separated by $\pm i\varphi_{a,b}$: \cref{eq_app_Q_full} is nonlocal and reduces to a Fokker--Planck equation only after truncation at small $\varphi_{a,b}$.

\subsection{Equations for the means}
We define $\mu_a=\int d^2\alpha\,d^2\beta\;\alpha Q=\langle\ha\rangle$ and $\mu_b=\langle\hb\rangle$. Integrating \cref{eq_app_Q_linear} by parts, $\int\alpha\,\partial_\alpha(A^\alpha Q)=-\langle A^\alpha\rangle$, while the other linear terms do not contribute to $\dot\mu_a$. For the nonlinear term, we use $\alpha f(D_L)Q=\pi^{-2}\bra{\alpha,\beta}f(\hat x)\hat\rho\,\ha\ket{\alpha,\beta}$, $\alpha f(D_R)Q=\pi^{-2}\bra{\alpha,\beta}\hat\rho f(\hat x)\ha\ket{\alpha,\beta}$, and the completeness relation $\int d^2\alpha\,d^2\beta\,\ket{\alpha,\beta}\bra{\alpha,\beta}=\pi^2$:
\begin{equation}
\begin{split}
    \int d^2\alpha\,d^2\beta\;\alpha\,[f(D_L)-f(D_R)]\,Q
    &=\langle[\ha,f(\hat x)]\rangle\\
    &=\varphi_a\langle f'(\hat x)\rangle .
\end{split}
\end{equation}
With $f=\sin$, we obtain the exact (Ehrenfest) equations
\begin{equation}
\begin{cases}
    \dot\mu_a &= -i\omega_a\mu_a+ig\varphi_a\bigl(\langle\cos\hat x\rangle-1\bigr),\\
    \dot\mu_b &= -\Bigl(i\omega_b+\frac{\kappa_b}{2}\Bigr)\mu_b-i\eps_d\cos(\omega_d t)\nonumber\\
    &\quad+ig\varphi_b\bigl(\langle\cos\hat x\rangle-1\bigr),
\end{cases}
\label{eq_app_mu}
\end{equation}
where the $-1$ comes from the linear term $g\hat x$ in $H_l$.

To close these equations, we assume a Gaussian state, characterized by the second central moments of $Q$,
\begin{equation}
\begin{aligned}
    \sigma_a^2&=\langle|\delta\alpha|^2\rangle, & \tilde\sigma_a^2&=\langle\delta\alpha^2\rangle, & c&=\langle\delta\alpha\,\delta\beta\rangle,\\
    \sigma_b^2&=\langle|\delta\beta|^2\rangle, & \tilde\sigma_b^2&=\langle\delta\beta^2\rangle, & d&=\langle\delta\alpha\,\delta\beta^*\rangle,
\end{aligned}
\end{equation}
with $\delta\alpha=\alpha-\mu_a$, $\delta\beta=\beta-\mu_b$. For a Gaussian state $\langle e^{i\hat x}\rangle=e^{i\bar X-K}$, with $\bar X=\langle\hat x\rangle=\varphi_a(\mu_a+\mu_a^*)+\varphi_b(\mu_b+\mu_b^*)$ and $K=\mathrm{Var}(\hat x)/2$. Because the moments of $Q$ are anti-normally ordered, $\mathrm{Var}(\ha+\ha\da)=2(\sigma_a^2+\mathrm{Re}\,\tilde\sigma_a^2)-1$, so that
\begin{equation}
\begin{split}
    K =\;& \varphi_a^2\Bigl(\sigma_a^2+\mathrm{Re}\,\tilde\sigma_a^2-\tfrac12\Bigr)
         +\varphi_b^2\Bigl(\sigma_b^2+\mathrm{Re}\,\tilde\sigma_b^2-\tfrac12\Bigr)\\
         &+2\varphi_a\varphi_b\,\mathrm{Re}(c+d),
\end{split}
\end{equation}
and $\langle\cos\hat x\rangle=e^{-K}\cos\bar X$.

The classical limit consists in neglecting the fluctuations of $\hat x$, $K\to0$, i.e.\ $\langle\cos\hat x\rangle\to\cos\bar X$. In terms of the moments of $Q$, this corresponds to an isotropic Gaussian with $\sigma_a^2=\sigma_b^2=1/2$ and $\tilde\sigma_a^2=\tilde\sigma_b^2=c=d=0$. Identifying $a\equiv\mu_a$ and $b\equiv\mu_b$, \cref{eq_app_mu} then reduce exactly to \cref{eq_classical_dynamics}. Note that a coherent state has $\sigma_{a,b}^2=1$, hence $K=(\varphi_a^2+\varphi_b^2)/2$: the first quantum correction renormalizes the cosine term by $e^{-(\varphi_a^2+\varphi_b^2)/2}$.

\section{Analytical derivation of the harmonic balance}
\label{app:derivation_harmonic_balance}

\paragraph*{Equations of motion.}
We consider the classical dynamics generated by
\begin{equation}
\begin{split}
    H(t) = {}& \omega_a |a|^2 + \omega_b |b|^2
    - 2E_J\sin\varepsilon_p\bigl[\sin\hat X - \hat X\bigr]\\
    &+ \varepsilon_d\cos(\omega_d t)\,(b+b^*),\\
    \hat X = {}& \varphi_a(a+a^*) + \varphi_b(b+b^*),
\end{split}
\label{eq:app_H}
\end{equation}
with $\dot a = -i\,\partial H/\partial a^*$ and $\dot b = -i\,\partial H/\partial b^* - \kappa_b b$.
We take $\omega_d=\omega_b=2\omega_a$. We rescale the fields, $\varphi_a a\to a$ and $\varphi_b b\to b$, use the dimensionless time $\tau=\omega_a t$, and move to the frame rotating at the mode frequencies, $a\to a\,e^{i\tau}$, $b\to b\,e^{2i\tau}$. With
\begin{equation}
\begin{aligned}
    \beta_a' &= \frac{2E_J\sin\varepsilon_p\,\varphi_a^2}{\omega_a},&\qquad
    \beta_b' &= \frac{2E_J\sin\varepsilon_p\,\varphi_b^2}{\omega_b},\\
    \varepsilon_d' &= \frac{\varepsilon_d\varphi_b}{\omega_b},&\qquad
    \kappa_b' &= \frac{\kappa_b}{\omega_b},
\end{aligned}
\end{equation}
the equations of motion read
\begin{equation}
\begin{split}
    \frac{da}{d\tau} &= i\beta_a'\,e^{i\tau}\bigl[\cos X - 1\bigr],\\
    \frac{db}{d\tau} &= 2i\beta_b'\,e^{2i\tau}\bigl[\cos X - 1\bigr]\\
    &\quad - i\varepsilon_d'\bigl(1+e^{4i\tau}\bigr) - 2\kappa_b' b,\\
    X(\tau) &= a\,e^{-i\tau} + b\,e^{-2i\tau} + \text{c.c.},
\end{split}
\label{eq:app_eom}
\end{equation}
where $X$ is the phase across the junction.

\paragraph*{Harmonic balance.}
We look for periodic solutions by inserting a finite Fourier Ansatz for $a(\tau)$ and $b(\tau)$ into \cref{eq:app_eom} and equating the Fourier components retained in the Ansatz. For any Ansatz, the junction phase takes the form
\begin{equation}
    X(\tau) = C + \sum_{j\ge1}\bigl(h_j e^{-ij\tau} + h_j^* e^{ij\tau}\bigr),
\end{equation}
with $C\in\mathbb R$ and $h_j = |h_j|e^{i\phi_j}$.
Using the Jacobi--Anger expansion $e^{iz\cos\vartheta}=\sum_{k} i^k J_k(z)e^{ik\vartheta}$ for each harmonic, one finds $\cos X=\sum_{N\in\mathbb Z}F_N e^{iN\tau}$ with
\begin{equation}
\begin{split}
    F_N = \sum_{\substack{\{k_j\}\in\mathbb Z\\ \sum_j j k_j = N}}
    &\Bigl[\prod_j J_{k_j}\bigl(2|h_j|\bigr)e^{-ik_j\phi_j}\Bigr]\\
    &\times\cos\Bigl(C + \frac{\pi}{2}\sum_j k_j\Bigr).
\end{split}
\label{eq:app_FN}
\end{equation}
Two properties will be used repeatedly. First, if the harmonics are phase locked as $h_j = r_j e^{ij\theta}$ with $r_j\in\mathbb R$ (signed), then $e^{-ik_j\phi_j}J_{k_j}(2|h_j|)=e^{-ijk_j\theta}J_{k_j}(2r_j)$ (if $r_j<0$, then $\phi_j=j\theta+\pi$ and we use $(-1)^kJ_k(z)=J_k(-z)$). The constraint $\sum_j jk_j=N$ then factors out the phase, $F_N = e^{-iN\theta}R_N$, with
\begin{equation}
    R_N = \sum_{\substack{\{k_j\}\in\mathbb Z\\ \sum_j j k_j = N}}
    \Bigl[\prod_j J_{k_j}(2r_j)\Bigr]
    \cos\Bigl(C + \frac{\pi}{2}\sum_j k_j\Bigr)\in\mathbb R .
    \label{eq:app_RN_general}
\end{equation} Second, at lowest order in the phase, $\cos X-1\simeq -X^2/2$, so that
\begin{equation}
    F_N \simeq \delta_{N,0} - \tfrac12\bigl(X^2\bigr)_N,
    \label{eq:app_FN_quad}
\end{equation}
where $(X^2)_N$ is the $e^{iN\tau}$ component of $X^2$.

\subsection{Cat states}
\label{app:cat}
\subsubsection{Derivation}
\label{app:cat_derivation}

The cat states are fixed points in the rotating frame, $a=a_0=|a_0|e^{i\theta_{a_0}}$ and $b=b_0$. Then $C=0$, $h_1=a_0$, $h_2=b_0$, and we keep only the constant component of \cref{eq:app_eom}, which selects $F_{-1}$ for the memory and $F_{-2}$ for the buffer:
\begin{equation}
\begin{cases}
    0 = F_{-1},\\
    0 = 2i\beta_b' F_{-2} - i\varepsilon_d' - 2\kappa_b' b_0,
\end{cases}
\end{equation}
with
\begin{equation}
\begin{split}
    F_N = \sum_{m\in\mathbb Z} &J_{N-2m}\bigl(2|a_0|\bigr)J_m\bigl(2|b_0|\bigr)\\
    &\times e^{-i(N-2m)\theta_{a_0}-im\theta_{b_0}}\cos\Bigl[\frac{(N-m)\pi}{2}\Bigr].
\end{split}
\end{equation}
We look for solutions with an empty buffer, $b_0=0$, so that only $m=0$ contributes. The memory equation is then satisfied identically, since $F_{-1}\propto\cos(-\pi/2)=0$, and $F_{-2}=-J_2(2|a_0|)e^{2i\theta_{a_0}}$. The buffer equation becomes
\begin{equation}
    J_2\bigl(2|a_0|\bigr)\,e^{2i\theta_{a_0}} = -\frac{\varepsilon_d'}{2\beta_b'}.
    \label{eq:app_cat_bessel}
\end{equation}
For $\varepsilon_d'>0$ this fixes $\theta_{a_0}=\pm\pi/2$ and $J_2(2|a_0|)=\varepsilon_d'/(2\beta_b')$, which has solutions only if $\varepsilon_d'/(2\beta_b')\le\max_z J_2(z)\simeq0.49$. Using $J_2(2x)=x^2/2-x^4/6+O(x^6)$, the lowest order gives the two cat states
\begin{equation}
\begin{split}
    &a_0 = \pm i\sqrt{\varepsilon_d'/\beta_b'},
    \quad\text{i.e.}\quad
    |\alpha|^2 = \frac{|a_0|^2}{\varphi_a^2} = \frac{\varepsilon_d}{2g_2},\\
    &g_2 = E_J\sin\varepsilon_p\,\varphi_a^2\varphi_b ,
\end{split}
\label{eq:app_cat}
\end{equation}
in the original units, and the next order gives $|a_0|^2=\tfrac32\bigl[1-\sqrt{1-4\varepsilon_d'/(3\beta_b')}\bigr]$.

\subsubsection{Equivalent Hamiltonian}
\label{app:cat_hamiltonian}

Expanding \cref{eq:app_H} to third order in $\hat X$ gives $-2E_J\sin\varepsilon_p[\sin\hat X-\hat X]\simeq (E_J\sin\varepsilon_p/3)\hat X^3$. For $\omega_b=2\omega_a$, the only term of $\hat X^3$ that is resonant in the rotating frame is $g_2(a^2b^*+\text{c.c.})$, which yields the two-photon Hamiltonian
\begin{equation}
\begin{split}
    H_\text{2ph} = {}& \omega_a|a|^2 + \omega_b|b|^2 + \varepsilon_d\cos(\omega_d t)(b+b^*)\\
    &+ g_2\bigl(a^2 b^* + a^{*2} b\bigr).
\end{split}
\end{equation}
With the same rescaling as above and after dropping the counter-rotating drive term $e^{4i\tau}$, its equations of motion read
\begin{equation}
\begin{cases}
    \dfrac{da}{d\tau} = -i\beta_a'\, a^* b,\\[2mm]
    \dfrac{db}{d\tau} = -i\beta_b'\, a^2 - i\varepsilon_d' - 2\kappa_b' b.
\end{cases}
\end{equation}
The fixed points satisfy $a_0^*b_0=0$ and $\beta_b'a_0^2+\varepsilon_d'-2i\kappa_b'b_0=0$, so that $b_0=0$ and $a_0^2=-\varepsilon_d'/\beta_b'$. This is exactly \cref{eq:app_cat}. In fact, since $J_2(2|a_0|)e^{2i\theta_{a_0}}=a_0^2/2+O(|a_0|^4)$, \cref{eq:app_cat_bessel} is the full-cosine version of the two-photon fixed-point condition: $H_\text{2ph}$ reproduces the cat states at leading order in the junction phase.

% =====================================================================
\subsection{Satellites}
\label{app:satellites}

\subsubsection{Derivation}
\label{app:sat_derivation}

\paragraph*{Ansatz.}
The satellites are periodic orbits of \cref{eq:app_eom} with period $2\pi$ (i.e.\ $2\pi/\omega_a$). We use
\begin{equation}
\begin{cases}
    a(\tau) = a_0 + a_1 e^{i\tau} + a_{-1}e^{-i\tau},\\
    b(\tau) = b_0 + b_{-1}e^{-i\tau} + b_2 e^{2i\tau}.
\end{cases}
\label{eq:app_sat_ansatz}
\end{equation}
In the laboratory frame:
\begin{itemize}
    \item $a_0$ and $b_0$ oscillate at $\omega_a$ and $\omega_b=2\omega_a$;
    \item $a_1$ and $b_2$ are static displacements;
    \item $a_{-1}$ is a memory component at $2\omega_a$;
    \item $b_{-1}$ is a buffer component at $3\omega_a$.
\end{itemize}

This choice follows from the nonlinearity. Through the quadratic term of $\cos X-1$, the leading harmonics $a_0$ and $b_0$ generate the frequencies $0,\omega_a,2\omega_a,3\omega_a,4\omega_a$. We keep the components detuned by at most $\omega_a$ from the mode frequencies, namely the static memory displacement, $a_{-1}$ and $b_{-1}$. We also keep the static buffer displacement $b_2$, which is driven by the same Fourier component $F_0$ as $a_1$ and contributes to the static junction phase.

The buffer component at $\omega_a$ is as close to resonance as these, but it is driven by $F_{-1}$, which vanishes on any solution (see below). An Ansatz containing it therefore returns a vanishing amplitude and misses the components above.

With \cref{eq:app_sat_ansatz}, the junction phase is
\begin{equation}
\begin{split}
    X(\tau) &= C + \bigl(a_0e^{-i\tau} + Y e^{-2i\tau} + b_{-1}e^{-3i\tau} + \text{c.c.}\bigr),\\
    C &= 2\,\mathrm{Re}(a_1+b_2),\qquad Y = b_0 + a_{-1}.
\end{split}
\end{equation}
The memory and buffer components at $2\omega_a$ add up in $X$, so that $h_1=a_0$, $h_2=Y$, $h_3=b_{-1}$, and \cref{eq:app_FN} becomes
\begin{equation}
\begin{split}
    F_N = \sum_{\substack{n,m,p\in\mathbb Z\\ n+2m+3p=N}}
    &J_n\bigl(2|a_0|\bigr)J_m\bigl(2|Y|\bigr)J_p\bigl(2|b_{-1}|\bigr)\\
    &\times e^{-i(n\theta_{a_0}+m\theta_{Y}+p\theta_{b_{-1}})}\\
    &\times\cos\Bigl[C+\frac{(n+m+p)\pi}{2}\Bigr].
\end{split}
\end{equation}

\paragraph*{Harmonic balance.}
We equate the components $e^{0}$, $e^{i\tau}$ and $e^{-i\tau}$ of the memory equation, and $e^{0}$, $e^{-i\tau}$ and $e^{2i\tau}$ of the buffer equation:
\begin{equation}
\begin{cases}
    0 = F_{-1},\\
    a_1 = \beta_a'\bigl(F_0-1\bigr),\\
    a_{-1} = -\beta_a' F_{-2},\\
    0 = 2i\beta_b'F_{-2} - i\varepsilon_d' - 2\kappa_b' b_0,\\
    \bigl(\kappa_b' - \tfrac{i}{2}\bigr)\,b_{-1} = i\beta_b' F_{-3},\\
    \bigl(1-i\kappa_b'\bigr)\,b_2 = \beta_b'\bigl(F_0-1\bigr).
\end{cases}
\label{eq:app_sat_hb}
\end{equation}

Three consequences follow directly.
\begin{itemize}
    \item Since $X$ is real, $F_0$ is real, and therefore $a_1\in\mathbb R$.
    \item The static components are slaved to each other, $b_2=\beta_b' a_1/[\beta_a'(1-i\kappa_b')]$, which gives
    \begin{equation}
        C = 2\gamma\,a_1,\qquad \gamma = 1+\frac{\beta_b'}{\beta_a'(1+\kappa_b'^2)} .
    \end{equation}
    \item The third and fourth lines of \cref{eq:app_sat_hb} share $F_{-2}$, so the drive fixes the memory component at $2\omega_a$ exactly, independently of how the cosine is truncated:
    \begin{equation}
        a_{-1} = -\frac{\beta_a'}{2\beta_b'}\bigl(\varepsilon_d' - 2i\kappa_b' b_0\bigr).
        \label{eq:app_am1}
    \end{equation}
\end{itemize}

\paragraph*{Phase reduction.}
We now use $\kappa_b'\ll1$. We set $\kappa_b'\to0$ in the last two lines of \cref{eq:app_sat_hb}, where it only gives $O(\kappa_b')$ corrections, and keep it in the drive equation, where it competes with $\varepsilon_d'$. We then look for phase-locked solutions
\begin{equation}
\begin{aligned}
    a_0 &= x\,e^{i\theta},&\qquad b_0 &= y\,e^{2i\theta},\\
    a_{-1} &= z\,e^{2i\theta},&\qquad b_{-1} &= w\,e^{3i\theta},
\end{aligned}
\end{equation}
with $x,y,z,w\in\mathbb R$,
so that $Y=(y+z)\,e^{2i\theta}\equiv \bar y\,e^{2i\theta}$. Signed amplitudes also cover relative phases shifted by $\pi$, which are then decided by the equations rather than fixed by hand.

With $F_N=e^{-iN\theta}R_N$, every equation becomes real except the drive equation, whose real and imaginary parts fix $\theta$:
\begin{equation}
\begin{cases}
    R_{-1} = 0,\\
    a_1 = \beta_a'\bigl(R_0-1\bigr),\\
    z = -\beta_a' R_{-2},\\
    2\beta_b' R_{-2} = \varepsilon_d'\cos2\theta,\\
    \varepsilon_d'\sin2\theta = -2\kappa_b' y,\\
    w = -2\beta_b' R_{-3},
\end{cases}
\label{eq:app_sat_bessel}
\end{equation}
with
\begin{equation}
\begin{split}
    R_N = \sum_{m,p\in\mathbb Z} &J_{N-2m-3p}(2x)\,J_m(2\bar y)\,J_p(2w)\\
    &\times\cos\Bigl[2\gamma a_1 + \frac{(N-m-2p)\pi}{2}\Bigr].
\end{split}
\label{eq:app_RN}
\end{equation}

Eliminating $\theta$ gives $2\beta_b'R_{-2}=\sigma\,\varepsilon_\text{eff}(y)$, with $\varepsilon_\text{eff}(y)=\sqrt{\varepsilon_d'^2-4\kappa_b'^2y^2}$ and $\sigma=\mathrm{sign}(\cos2\theta)$. This requires $2\kappa_b'|y|\le\varepsilon_d'$, and it gives
\begin{equation}
    z = -\sigma\,\frac{\beta_a'}{2\beta_b'}\,\varepsilon_\text{eff}(y),
    \qquad y = \bar y - z .
    \label{eq:app_z}
\end{equation}

The two signs correspond to two families of periodic orbits:
\begin{itemize}
    \item $\sigma=-1$, $\theta\simeq\pm\pi/2$: the cat states. Setting $a_{\pm1}=b_2=b_{-1}=0$ recovers \cref{app:cat_derivation}.
    \item $\sigma=+1$, $\theta\simeq0,\pi$: the satellites. They are rotated by $90^\circ$ with respect to the cats, up to the dissipative tilt $\theta_{b_0}=2\theta=-\arcsin(2\kappa_b' y/\varepsilon_d')$, with $\theta_{a_0}=\theta_{b_0}/2 \ [\pi]$.
\end{itemize}
\Cref{eq:app_sat_bessel} is invariant under $(x,w)\to(-x,-w)$, i.e.\ under $a\to-a$, so the satellites come in pairs related by the parity of the memory.

\paragraph*{Leading-order solution.}
At lowest order in the junction phase, \cref{eq:app_FN_quad} gives
\begin{equation}
\begin{split}
    R_{-1}&\simeq -x(C+\bar y)-\bar y w,\\
    R_0-1&\simeq -\Bigl(\frac{C^2}{2}+x^2+\bar y^2+w^2\Bigr),\\
    R_{-2}&\simeq -\Bigl(C\bar y+\frac{x^2}{2}+xw\Bigr),\\
    R_{-3}&\simeq -(Cw+x\bar y).
\end{split}
\end{equation}
For $\kappa_b'\to0$ in $\gamma$, we have $\gamma\beta_a'=\beta_a'+\beta_b'\equiv\beta_\Sigma$, and \cref{eq:app_sat_bessel} becomes the polynomial system
\begin{equation}
\begin{cases}
    x\,(C+\bar y) + \bar y\,w = 0,\\
    C = -2\beta_\Sigma\Bigl(\dfrac{C^2}{2}+x^2+\bar y^2+w^2\Bigr),\\
    -2\beta_b'\Bigl(C\bar y+\dfrac{x^2}{2}+xw\Bigr) = \sigma\,\varepsilon_\text{eff}(y),\\
    w = 2\beta_b'\,(Cw+x\bar y),
\end{cases}
\label{eq:app_sat_poly}
\end{equation}
completed by \cref{eq:app_z}.

For $x\neq0$, let $u=w/x$ and $\eta=2\beta_b'\bar y$. The first and last lines give $C=-\bar y(1+u)$ and $\eta u^2+(1+\eta)u-\eta=0$, whose root continuously connected to $u=0$ is
\begin{equation}
    u(\bar y) = \frac{\sqrt{(1+\eta)^2+4\eta^2}-(1+\eta)}{2\eta} = \eta + O(\eta^2).
\end{equation}
The third line with $\sigma=+1$ then fixes the memory amplitude,
\begin{equation}
    x^2 = \frac{2(1+u)\,\bar y^2 - \varepsilon_\text{eff}(y)/\beta_b'}{1+2u},
\end{equation}
and the second line reduces the problem to a single scalar equation for the total $2\omega_a$ amplitude $\bar y$:
\begin{equation}
    \bar y\,(1+u) = 2\beta_\Sigma\Bigl[\Bigl(1+\frac{(1+u)^2}{2}\Bigr)\bar y^2 + \bigl(1+u^2\bigr)\,x^2\Bigr].
    \label{eq:app_sat_scalar}
\end{equation}
Here $y$ is related to $\bar y$ by $y=\bar y+\beta_a'\varepsilon_\text{eff}(y)/(2\beta_b')$ from \cref{eq:app_z}. Since $C<0$ by the second line of \cref{eq:app_sat_poly}, the physical root has $\bar y>0$ and must satisfy $x^2>0$.

All other amplitudes follow:
\begin{equation}
\begin{aligned}
    a_0 &= \pm x\,e^{i\theta}, & b_0 &= y\,e^{2i\theta},\\
    b_{-1} &= \pm u\,x\,e^{3i\theta}, & a_{-1} &= -\frac{\beta_a'\varepsilon_\text{eff}(y)}{2\beta_b'}\,e^{2i\theta},\\
    a_1 &= -\frac{(1+u)\,\bar y}{2\gamma}, & b_2 &= \frac{\beta_b'}{\beta_a'}\,a_1 .
\end{aligned}
\label{eq:hb_analytic_complex}
\end{equation}
To lowest order in $\eta$, $|b_{-1}|\simeq2\beta_b'|a_0|\bar y$.

In the limit $u\to0$ and $\kappa_b'\to0$, \cref{eq:app_sat_scalar} reduces to $7\beta_\Sigma \bar y^2-\bar y-2\beta_\Sigma\varepsilon_d'/\beta_b'=0$, i.e.
\begin{equation}
\begin{split}
    \bar y &= \frac{1+\sqrt{1+56\,\beta_\Sigma^2\,\varepsilon_d'/\beta_b'}}{14\,\beta_\Sigma},\\
    |a_0|^2 &= 2\bar y^2 - \frac{\varepsilon_d'}{\beta_b'},
    \qquad |b_0| = \bar y + \frac{\beta_a'\varepsilon_d'}{2\beta_b'} .
\end{split}
\label{eq:hb_analytic_simple}
\end{equation}

\paragraph*{Orbits and validity.}
In the rotating frame, the memory follows $a_0+a_1e^{i\tau}+a_{-1}e^{-i\tau}$, an ellipse centered on $a_0$ with semi-axes $|a_1|\pm|a_{-1}|$. The buffer follows $b_0+b_2e^{2i\tau}+b_{-1}e^{-i\tau}$, which describes a trefoil, matching \cref{fig_limit_cycles} .

Since $\bar y\sim1/(7\beta_\Sigma)$, the expansion in the junction phase is quantitatively controlled only when $|a_0|,\bar y,|C|\ll1$. Otherwise, the full Bessel system \cref{eq:app_sat_bessel} has to be solved numerically, using the solution of \cref{eq:app_sat_scalar} as the initial guess.

\subsubsection{Equivalent Hamiltonian}
\label{app:sat_hamiltonian}

We look for a truncated Hamiltonian whose harmonic balance with the Ansatz \cref{eq:app_sat_ansatz} gives exactly the approximated equations, i.e.\ \cref{eq:app_sat_hb} with $F_N$ replaced by \cref{eq:app_FN_quad}. This is the case for the cubic expansion of \cref{eq:app_H}, $-2E_J\sin\varepsilon_p[\sin\hat X-\hat X]\to(E_J\sin\varepsilon_p/3)\hat X^3$, because its equations of motion are \cref{eq:app_eom} with $\cos X-1\to-X^2/2$. In normal-ordered monomials it reads
\begin{equation}
\begin{split}
    H_\text{trunc} = {}& \omega_a|a|^2 + \omega_b|b|^2 + \varepsilon_d\cos(\omega_d t)(b+b^*)\\
    &+ E_J\sin\varepsilon_p\Bigl[
    \varphi_a^2\varphi_b\,a^2 b^*
    + \varphi_a^3\,a^2a^*
    + \varphi_b^3\,b\,b^{*2}\\
    &\qquad + \varphi_a\varphi_b^2\,a\,b^{*2}
    + 2\varphi_a^2\varphi_b\,|a|^2 b\\
    &\qquad + 2\varphi_a\varphi_b^2\,a\,|b|^2
    + \text{c.c.}\Bigr]\\
    &+ E_J\sin\varepsilon_p\Bigl[
    \frac{\varphi_a^3}{3}\,a^3
    + \varphi_a^2\varphi_b\,a^2 b
    + \varphi_a\varphi_b^2\,a\,b^2\\
    &\qquad + \frac{\varphi_b^3}{3}\,b^3
    + \text{c.c.}\Bigr].
\end{split}
\label{eq:app_Htrunc}
\end{equation}

Inserting \cref{eq:app_sat_ansatz} and projecting as in \cref{eq:app_sat_hb}, its harmonic balance reads, before phase reduction,
\begin{equation}
\begin{split}
    0 &= C a_0 + Y a_0^* + Y^* b_{-1},\\
    a_1 &= -\beta_a'\Bigl(\frac{C^2}{2} + |a_0|^2 + |Y|^2 + |b_{-1}|^2\Bigr),\\
    a_{-1} &= \beta_a'\Bigl(CY + \frac{a_0^2}{2} + a_0^* b_{-1}\Bigr),\\
    0 &= -i\beta_b'\Bigl(CY + \frac{a_0^2}{2} + a_0^* b_{-1}\Bigr)\\
    &\quad - \frac{i\varepsilon_d'}{2} - \kappa_b' b_0,\\
    \bigl(\kappa_b' - \tfrac{i}{2}\bigr)\,b_{-1} &= -i\beta_b'\bigl(C b_{-1} + a_0 Y\bigr),\\
    \bigl(1-i\kappa_b'\bigr)\,b_2 &= -\beta_b'\Bigl(\frac{C^2}{2} + |a_0|^2 + |Y|^2 + |b_{-1}|^2\Bigr),
\end{split}
\label{eq:app_sat_hb_quad}
\end{equation}
with $C=2\mathrm{Re}(a_1+b_2)$ and $Y=b_0+a_{-1}$. After phase locking, these equations are exactly \cref{eq:app_sat_poly,eq:app_z}.

The two groups of terms in \cref{eq:app_Htrunc} play different roles.
\begin{itemize}
    \item The first bracket contains every term that contributes to the equations of the oscillating amplitudes $a_0$, $a_{-1}$, $b_0$ and $b_{-1}$:
    \begin{itemize}
        \item the two-photon exchange $a^2b^*$, which alone gives the cat states of \cref{app:cat_hamiltonian};
        \item together with $a|b|^2$, the same term produces the three-wave mixing $a_0Y\leftrightarrow b_{-1}$ responsible for the sideband at $3\omega_a$;
        \item evaluated on the static displacements, all the terms of this bracket produce the frequency shifts $Ca_0$, $CY$ and $Cb_{-1}$.
    \end{itemize}
    \item The second bracket, $a^3$, $a^2b$, $ab^2$, $b^3$ and their complex conjugates, contributes only to the two static equations (for $a_1$ and $b_2$), through products of $a_1$ and $b_2$. These terms are nevertheless required: together with the static parts of the first bracket, they build the complete $C^2/2$ term, which is of the same order as the other terms.
\end{itemize}

\begin{figure}[tbp]
    \centering
    \includegraphics[width=\columnwidth]{Chaos/figures_article/classical_limit_cycles_analytic_HB_12_10_2026.pdf}
    \caption{Prediction of the harmonic balance. In yellow the numerical solution with many harmonics, in green the the numerical solution of the system \cref{eq:hb_analytic_complex} and in pink the analytical solution \cref{eq:hb_analytic_simple}. Despite its simplicity, we get a good estimate of the position and size of the satellite limit cycles for the memory.}
    \label{fig:hb_analytic_limit_cycles}
\end{figure}

\section{Period-doubled satellite limit-cycle}

\begin{figure}[tbp]
    \centering
    \includegraphics[width=\columnwidth]{Chaos/figures_article/classical_limit_cycles_period_doubled_12_10_2026.pdf}
    \caption{Example of limit-cycles after the period-doubling}
    \label{fig:limit_cycles_period_doubled}
\end{figure}

\section{Lyapunov heatmap with initial points at the center of the cat lobes}

We check that the fringes are indeed the most fragile feature of the cat qubit we work with. To show this, we compute the maximal Lyapunov exponent when the trajectory is started close to a cat limit-cycle, instead of starting at the origin of phase space like in the main text. As shown on \cref{fig:lyapunov_heatmap_cat_start}, the threshold for chaotic dynamics is higher, meaning we can work with higher values of the two-photon exchange rate $g_2$ or buffer drive amplitude $\eps_d$ and observe regular trajectories.

\TMD{Other reason why it is important to check this?}

\begin{figure}[tbp]
    \centering
    \includegraphics[width=\columnwidth]{Chaos/figures_article/heatmap_lyapunov_cat_start_restyled_12_10_2026.pdf}
    \caption{Lyapunov heatmap when the trajectory is initialized close to a cat limit-cycle}
    \label{fig:lyapunov_heatmap_cat_start}
\end{figure}

\section{Lyapunov heatmap where the instability matches the blue to red transition}

\TMD{Check if satellites do not exist in the light blue region, with less time for the computation of the Lyapunov exponent.}

\iffalse
\subsection{Method}
\label{app:lyap_method}

\paragraph*{Equations of motion.}
We integrate the classical equations of the two modes with the full Josephson nonlinearity. With the dimensionless quadratures $x_a$ and $x_b$, related to the complex amplitudes by $a = (x_a + i\dot x_a/\omega_a)/\sqrt2$ and $b = (x_b + i\dot x_b/\omega_b)/\sqrt2$, they read
\begin{equation}
\begin{cases}
    \ddot x_a = -\omega_a^2 x_a + 2\sqrt2\,E_J\sin\varepsilon_p\,\omega_a\varphi_a\,\bigl[\cos X - 1\bigr],\\
    \ddot x_b = -\omega_b^2 x_b - \sqrt2\,\omega_b\,\varepsilon_d\cos(\omega_d t) + 2\sqrt2\,E_J\sin\varepsilon_p\,\omega_b\varphi_b\,\bigl[\cos X - 1\bigr] - 2\kappa_b\dot x_b,
\end{cases}
\label{eq:app_lyap_eom}
\end{equation}
with junction phase $X = \sqrt2(\varphi_a x_a + \varphi_b x_b)$. The pump frequency is $\omega_p = 2\omega_a - \omega_d = 0$ and the drive frequency is $\omega_d = \omega_b = 2\omega_a$. The memory is undamped ($\kappa_a = 0$). The buffer damping corresponds to $\dot b = \dots - \kappa_b b$.

The parameters are $\omega_a/2\pi = 4.033~\mathrm{GHz}$, $\varphi_a = 0.11$, $\varphi_b = 0.204$, $E_J/2\pi = 55.68~\mathrm{GHz}$ and $\kappa_b = \pi/10.4~\mathrm{ns}^{-1}$ ($\kappa_b/2\pi \simeq 48~\mathrm{MHz}$). The maps are drawn against
\begin{equation}
    g_2 = J_1(\varepsilon_p)\,E_J\varphi_a^2\varphi_b, \qquad \frac{\varepsilon_d}{2g_2},
\end{equation}
on a grid of $91\times86$ points: $\varepsilon_p \in [0.2, 1.1]$ in steps of $0.01$ ($g_2/2\pi$ from $13.7$ to $64.7~\mathrm{MHz}$) and $\varepsilon_d/g_2 \in [5, 13.5]$ in steps of $0.1$.

\paragraph*{Largest Lyapunov exponent.}
For each grid point we estimate the largest Lyapunov exponent $\lambda$ of \cref{eq:app_lyap_eom} with the two-trajectory (Benettin) method, as implemented in \texttt{ChaosTools.jl}:
\begin{enumerate}
    \item The trajectory is first evolved over a transient $T_\text{tr} = 15/\kappa_b \simeq 50~\mathrm{ns}$, which is discarded.
    \item A test trajectory is started at a distance $d_0 = 10^{-7}$ from the reference, in the Euclidean norm of the state $(x_a, x_b, \dot x_a, \dot x_b)$.
    \item Both are integrated together over $T = 30/\kappa_b \simeq 99~\mathrm{ns}$, about $400$ memory periods. Their distance $d$ is checked every integration step, $\Delta t = 1/(34\sqrt2\,\omega_a) \simeq 0.8~\mathrm{ps}$.
    \item Whenever $d$ exceeds $d_\text{up} = 10^{-4}$, the test trajectory is brought back along the separation to distance $d_0$, and $\ln(d/d_0)$ is accumulated.
\end{enumerate}
The exponent is the accumulated growth rate,
\begin{equation}
    \lambda = \frac{1}{T}\sum_i \ln\frac{d_i}{d_0}.
\end{equation}
The maps show $\mathrm{sign}(\lambda)$: positive for chaotic trajectories, negative for trajectories converging to a fixed point or limit cycle of the stroboscopic dynamics. This is a finite-time estimate over a window of $45/\kappa_b$ including the transient. A trajectory that is still chaotic at the end of the window is counted as chaotic, even if it would eventually settle.

\paragraph*{Initial conditions.}
We compute three maps, which differ only by the initial state of every trajectory.
\begin{enumerate}
    \item \emph{Vacuum:} $a(0) = b(0) = 0$.
    \item \emph{Cat:} one of the two cat attractors, $a(0) = i|\alpha|$, $b(0) = 0$. Its amplitude follows from the cat solution of the harmonic balance (\cref{app:cat_derivation}), $\beta_b'|a_0|^2 = \varepsilon_d'$, which in the units of \cref{eq:app_lyap_eom} reads
    \begin{equation}
        |\alpha|^2 = \frac{\varepsilon_d}{2g_2}\,\frac{J_1(\varepsilon_p)}{\sin\varepsilon_p}.
    \end{equation}
    The ratio $J_1(\varepsilon_p)/\sin\varepsilon_p$ arises because $g_2$ is defined with $J_1(\varepsilon_p)$, while for $\omega_p = 0$ the pump enters \cref{eq:app_lyap_eom} through $\sin\varepsilon_p$. It varies between $0.50$ and $0.53$ over the grid.
    \item \emph{Satellite:} the point $\tau = 0$ of the satellite orbit of \cref{app:satellites}, solved at every grid point from the full system (finite $u$, with the dissipative tilt $2\theta = -\arcsin(2\kappa_b'y/\varepsilon_d')$):
    \begin{equation}
        a(0) = \frac{a_0 + a_1 + a_{-1}}{\varphi_a}, \qquad b(0) = \frac{b_0 + b_{-1} + b_2}{\varphi_b}.
    \end{equation}
    The analytic amplitudes lie about $10\%$ below those of the simulated orbit, so the trajectory starts close to the satellite and relaxes onto it if the satellite is stable. On $21.5\%$ of the grid, at the lowest $g_2$, the equations have no admissible root ($x^2 \le 0$ or $2\kappa_b'|y| > \varepsilon_d'$). No satellite exists there, and these points are left blank.
\end{enumerate}

\paragraph*{End state.}
For the satellite map, a trajectory with $\lambda < 0$ may still end on a cat rather than on the satellite. Each trajectory is therefore integrated again up to $T_\text{tr} + T$, and $\langle |b| \rangle$ is averaged over the last memory period. This quantity is $\simeq 0.1$ on a cat and $1.6$--$2.1$ on a satellite. Each point is classified as
\begin{itemize}
    \item \emph{chaotic} if $\lambda > 0$;
    \item \emph{satellite} if $\lambda < 0$ and $\langle|b|\rangle > 0.8$;
    \item \emph{cat} otherwise.
\end{itemize}

\subsection{Results}
\label{app:lyap_results}

\paragraph*{Vacuum start.}
Starting from $a = b = 0$, $16.5\%$ of the grid is chaotic. The chaotic region is bounded from below by an onset line, the smallest chaotic $\varepsilon_d/2g_2$ at each $\varepsilon_p$. This line runs from $\varepsilon_d/2g_2 = 6.75$ at $g_2/2\pi \simeq 44~\mathrm{MHz}$ down to $\varepsilon_d/2g_2 \simeq 3.4$ at the largest $g_2$. Below this line every trajectory reaches a regular attractor.

\paragraph*{Cat start.}
Starting on a cat, the chaotic fraction falls to $3.4\%$. Of the points that are chaotic from the vacuum, $1077$ become regular, while only $55$ points are chaotic from the cat but regular from the vacuum. The remaining chaos is confined to $g_2/2\pi \gtrsim 55~\mathrm{MHz}$ and $\varepsilon_d/2g_2 \gtrsim 5.3$.

The cat attractor therefore remains stable over most of the region where the vacuum start is chaotic. There, chaos is a property of the initial condition, not an instability of the cat. The dynamics is multistable: a stable cat coexists with a chaotic set whose basin contains the vacuum.

\paragraph*{Satellite start.}
\fi
\Cref{fig:app_lyap_satellite} shows the map for trajectories started on the satellite. Of the points where a satellite solution exists, $30.8\%$ end on the satellite, $29.2\%$ on a cat, and $40.0\%$ are chaotic. From low to high $g_2$, the map shows four regions.
\begin{itemize}
    \item A band about $8~\mathrm{MHz}$ wide along the edge of the blank region, where the trajectory relaxes to a cat. There the analytic satellite exists but the actual satellite is unstable or absent, so the starting point lies in the basin of a cat.
    \item A band where the trajectory stays on the satellite ($\lambda < 0$, large $\langle|b|\rangle$), up to $g_2/2\pi \simeq 41~\mathrm{MHz}$ at $\varepsilon_d/2g_2 = 6.75$ and $\simeq 45~\mathrm{MHz}$ at $\varepsilon_d/2g_2 = 2.5$. Its upper edge coincides, within one or two grid cells, with the instability line of the satellite branch obtained independently from Hill's method (\cref{app:hill}). Along that branch, the satellite first period-doubles. At $\varepsilon_d = 10.12\,g_2$ a Floquet multiplier crosses $-1$ at $g_2/2\pi \simeq 41.0~\mathrm{MHz}$. The period-two orbit then loses stability through a complex pair of multipliers leaving the unit circle (a Neimark--Sacker bifurcation) at $g_2/2\pi \simeq 41.7~\mathrm{MHz}$. Just beyond that line, the trajectories become chaotic.
    \item A region of uniform chaos, up to $g_2/2\pi \simeq 50~\mathrm{MHz}$.
    \item At larger $g_2$, a fine-grained mixture of chaotic, cat and satellite end states between neighbouring grid points.
\end{itemize}
The agreement between the onset of chaos and Hill's instability line identifies the destabilisation of the satellite branch as the route to chaos. The vacuum-start onset lies at larger $g_2$ than the satellite instability. This is consistent with the vacuum having to fall into the basin of the chaotic set first, rather than starting on the branch that generates it.

\paragraph*{Limitations.}
The fine-grained region at large $g_2$ should not be read as a sharp phase diagram. There the finite-time exponent is evaluated over $45/\kappa_b$, which is comparable to the lifetime of transient chaos at these drive strengths. For example, $\simeq 42/\kappa_b$ at $\varepsilon_p = 0.802$, $g_2/2\pi \simeq 50.8~\mathrm{MHz}$ (\cref{sec:transient_entropy}). Neighbouring points can then fall on either side of the end of the transient. In the same way, the cat and satellite basins are probed with a single initial condition per grid point. The maps thus show which attractor that particular initial condition reaches within the window, not the full basin structure.

\begin{figure}[tbp]
    \centering
    \includegraphics[width=\columnwidth]{Chaos/figures_article/heatmap_lyapunov_satellite_start_12_10_2026.pdf}
    \caption{Trajectories started on the satellite. Left: sign of the largest Lyapunov exponent (blue $\lambda<0$, red $\lambda>0$). Right: end state (satellite, cat or chaotic). Blank: no satellite solution. Black line: instability of the satellite branch from Hill's method (\cref{app:hill}).}
    \label{fig:app_lyap_satellite}
\end{figure}


\section{Computation of IPR as a signature of chaos}
\label{app:ipr}

\subsection{Definition}

For a pure state $|\psi\rangle$ expanded on an orthonormal basis
$\{|k\rangle\}$, the inverse participation ratio (IPR) is
\begin{equation}
  \mathrm{IPR} = \sum_k |\langle k|\psi\rangle|^4 .
  \label{eq:ipr_pure}
\end{equation}
It equals $1$ when $|\psi\rangle$ is a single basis state and $1/N$ when
$|\psi\rangle$ is spread evenly over $N$ of them. Its inverse,
$N_\mathrm{eff} = 1/\mathrm{IPR}$, counts the basis states that effectively
take part in $|\psi\rangle$. The memory is entangled with the lossy buffer, so
its reduced state $\hat\rho_a = \mathrm{Tr}_b\,\hat\rho$ is mixed. We
therefore use the natural generalization
\begin{equation}
  \mathrm{IPR}(t) = \sum_{k=0}^{n_a-1} p_k(t)^2 ,
  \qquad
  p_k(t) = \langle k|\hat\rho_a(t)|k\rangle ,
  \label{eq:ipr_mixed}
\end{equation}
where $|k\rangle$ are the Fock states of the memory. For a pure state,
Eq.~\eqref{eq:ipr_mixed} reduces to Eq.~\eqref{eq:ipr_pure}. Only the Fock
populations enter: Eq.~\eqref{eq:ipr_mixed} measures how far the photon-number
distribution of the memory is spread, independently of the coherences
between Fock states.

\subsection{IPR as a probe of quantum chaos}

The participation ratio was introduced to quantify the localization of
vibrational modes~\cite{Bell1970} and electronic wave
functions~\cite{Wegner1980,Kramer1993}. It became a standard measure of
eigenstate delocalization in quantum chaos~\cite{Izrailev1990,Haake2010}. In a
regular (integrable) system, eigenstates and wave packets stay localized in a
suitable unperturbed basis, and the IPR remains of order one. In a chaotic
system they spread ergodically over the energetically available states, and
the IPR decreases towards its random-state value, $2/(N+1)$ for a random
complex vector in dimension $N$~\cite{Izrailev1990,Santos2010,Borgonovi2016}.
Applied to a time-evolved state instead of eigenstates, a decay of
$\mathrm{IPR}(t)$ quantifies the dynamical spreading of an initially localized
state over the basis~\cite{Borgonovi2016}. It complements the
out-of-time-order correlator, which probes operator growth.

For a cat qubit, the Fock basis of the memory is the natural reference. In the
stabilized regime, the two-photon dissipation confines the memory to the cat
manifold spanned by $|\pm\alpha\rangle$. Its photon-number distribution
remains that of the initial cat, a Poisson distribution of mean $|\alpha|^2$
restricted to one parity. A spreading of the Fock distribution beyond it,
i.e.\ a decrease of the IPR, signals that the dynamics leaves the cat
manifold.

\subsection{Results}

Figure~\ref{fig:ipr} shows $\mathrm{IPR}(T)$ over the
$(\varepsilon_p, \varepsilon_d/g_2)$ plane. Most of the vertical gradient
comes from the initial state: a larger drive prepares a larger cat, whose
Fock distribution is wider, and $\mathrm{IPR}(0)$ decreases from $0.526$ at
$r = 2$ to $0.270$ at $r = 9$. The dynamics is read from the departure of
$\mathrm{IPR}(T)$ from $\mathrm{IPR}(0)$ along each row.

At weak pump amplitude the Fock distribution is frozen. At
$\varepsilon_p = 0.2$, $\mathrm{IPR}(T)/\mathrm{IPR}(0) \geq 0.96$ for every
drive strength. This is the expected behavior of a stabilized cat qubit,
whose state remains in the cat manifold.

As $\varepsilon_p$ increases, $\mathrm{IPR}(T)$ decreases for every
$r \geq 4$, and the decrease is larger at stronger drive. It proceeds in two
stages. Up to $\varepsilon_p \approx 0.8$, $\mathrm{IPR}(T)$ decreases smoothly,
by at most $\sim 25\%$ relative to its value at $\varepsilon_p = 0.2$. Above
$\varepsilon_p \approx 0.85$ it drops faster. At $\varepsilon_p = 1.08$ it
reaches $\mathrm{IPR}(T)/\mathrm{IPR}(0) = 0.74$ for $r = 4$ and $0.38$ for
$r = 9$. In that corner, the memory effectively occupies
$N_\mathrm{eff} = 1/\mathrm{IPR} \approx 9.7$ Fock states, against $3.7$ for
the initial cat. This is the strong spreading of the photon-number
distribution expected from chaotic dynamics, which mixes the cat with Fock
states outside its parity sector and its Poisson envelope. The small
non-monotonic features along the rows, e.g.\ near $\varepsilon_p \approx 0.93$,
are expected from a snapshot at a single time of a state that is still
evolving.

At the smallest cat sizes, the behavior is different. For $r = 2$
($|\alpha|^2 = 1$), $\mathrm{IPR}(T)$ instead increases with $\varepsilon_p$,
by up to $7\%$: the Fock distribution concentrates rather than spreads. For
$r = 3$, the IPR shows only a shallow minimum, $8\%$ below
$\mathrm{IPR}(0)$, around $\varepsilon_p \approx 0.8$. The spreading
therefore requires both a large pump amplitude and a sufficiently large
cat.

The same simulations give the Wehrl entropy $S_W(T)$ of the memory. Over the
$184$ points of the grid, $S_W(T)$ and $-\ln\mathrm{IPR}(T)$ are strongly
correlated (Pearson coefficient $0.98$). Delocalization in the Fock basis
and spreading of the Husimi distribution in phase space thus tell the same
story.

Two caveats apply. First, the IPR is basis dependent: a decrease signals
spreading in the Fock basis of the memory, which can also come from a change
of the mean photon number or from the entanglement with, and the loss
through, the buffer. Its interpretation as a chaos indicator rests on the
comparison with the stabilized regime, where these effects are present but
leave the IPR unchanged. Second, $T = 0.2/\kappa_b$ is short compared with
the buffer lifetime. Figure~\ref{fig:ipr} therefore probes the onset of the
spreading, not its long-time saturation.

\begin{figure}[tbp]
  \centering
  \includegraphics[width=\columnwidth]{Chaos/figures_article/IPR_heatmap_12_10_2026.pdf}
  \caption{Inverse participation ratio
    $\mathrm{IPR}(T) = \sum_k \langle k|\hat\rho_a(T)|k\rangle^2$ of the
    memory in its Fock basis, at $T = 0.2/\kappa_b$, as a function of the pump
    amplitude $\varepsilon_p$ and the drive strength $\varepsilon_d/g_2$.
    The memory starts in the even cat state of amplitude
    $\alpha = \sqrt{\varepsilon_d/2g_2}$. Truncations $n_a = 29$, $n_b = 17$.}
  \label{fig:ipr}
\end{figure}


\section{OTOC but computed on a coherent state and not a cat}
\label{app:OTOC_coherent}

One could assert that the computation of the OTOC is dependent on the state it gets averaged on. Furthermore, we showed the results for a cat state, a bimodally distributed state and could be different for a more standard unimodally distributed state such as a coherent state. In this appendix, we present on \cref{fig:OTOC_coherent} the maximal Lyapunov exponent derived from the OTOC when averaged on the coherent state $\ket{\alpha}$, building block of the cat state. We see no major changes.

\begin{figure}[tbp]
    \centering
    \includegraphics[width=\columnwidth]{Chaos/figures_article/Lyapunov_from_OTOC_init_coherent_12_10_2026.pdf}
    \caption{OTOC averaged on a coherent state (lobe of the cat state)}
    \label{fig:OTOC_coherent}
\end{figure}

\section{Derivation of the parity-protected Hamiltonian}
\label{app:parity_circuit}

In this appendix we derive the Hamiltonian of the mirror-symmetric
circuit whose nonlinearity is even in the memory quadratures,
\begin{equation}
  -4E_J \sin[\epsilon(t)]\cos(\varphi_a \hat{x}_a)\sin(\varphi_b \hat{x}_b),
  \label{eq:app_target}
\end{equation}
by combining two identical asymmetrically threaded SQUIDs (ATS).
The circuit is treated classically and quantized only at the end
(Sec.~\ref{app:quantization}). We use the reduced flux quantum
$\phi_0=\hbar/2e$, write phases as $\phi=\Phi/\phi_0$, and define
$E_L=\phi_0^2/L$ and $E_C=e^2/2C$.

% --- Optional figure -------------------------------------------------
% \begin{figure}[t]
%   \centering
%   \includegraphics[width=0.8\columnwidth]{fig_two_ATS_circuit}
%   \caption{Mirror-symmetric circuit (configuration A). Two identical
%   ATS elements connect islands 1 and 2 to ground and receive identical
%   external fluxes. They are coupled by $C_c$ (optionally $L_c$), and the
%   transmission line couples symmetrically to both inductor nodes through
%   $C_p$. The memory is the antisymmetric mode, the buffer the symmetric one.}
%   \label{fig:app_circuit}
% \end{figure}

\subsection{ATS element}

Each element $i=1,2$ consists of two junctions ($E_J$, capacitance $C_J$)
and a central inductor $L$ in parallel between an island and ground, with
sub-nodes $\Phi_{i,1}$, $\Phi_{i,2}$ (inductor) and $\Phi_{i,3}$. The loop
fluxes impose $\Phi_{i,2}-\Phi_{i,1}=\Phi_{l,i}(t)$ and
$\Phi_{i,3}-\Phi_{i,2}=\Phi_{r,i}(t)$. Defining
$\Phi_{\Sigma,i}=(\Phi_{l,i}+\Phi_{r,i})/2$,
$\Phi_{\Delta,i}=(\Phi_{r,i}-\Phi_{l,i})/2$ and the island coordinate
$\Phi_i=(\Phi_{i,1}+\Phi_{i,3})/2$, one has
$\Phi_{i,1,3}=\Phi_i\mp\Phi_{\Sigma,i}$ and $\Phi_{i,2}=\Phi_i-\Phi_{\Delta,i}$.
Assuming equal junction capacitances and that the external circuit
attaches to the inductor sub-node $\Phi_{i,2}$, the element Lagrangian is
$\mathcal{L}_{\mathrm{ATS},i}=C_J\dot{\Phi}_i^2-U_i(\phi_i,t)$, up to a
term depending only on the external fluxes, with
\begin{equation}
  U_i(\phi_i,t)=\frac{E_L}{2}\left(\phi_i-\phi_{\Delta,i}\right)^2
  -2E_J\cos\phi_{\Sigma,i}(t)\,\cos\phi_i .
  \label{eq:app_U}
\end{equation}
The $\Sigma$-flux generates no charge drive on $\Phi_i$, whereas the
$\Delta$-flux enters both the inductive potential and, because the
external circuit is attached at $\Phi_i-\Phi_{\Delta,i}(t)$, the kinetic
terms of the coupling elements. These terms are kept explicitly, which is
equivalent to the irrotational construction~\cite{You2019May}. Unequal
junction capacitances only add identical offsets to $\phi_{\Sigma,i}$ in
both elements, which affect the symmetric coordinate alone and leave the
parity argument unchanged.

\subsection{Circuit Lagrangian and bias point}

The inductor nodes carry a shunt capacitance $C_s$, are coupled to each
other by $C_c$ (and optionally $L_c$), and each couples through $C_p$ to
the port node at voltage $V_p(t)$. With identical $\Delta$-fluxes
$\Phi_{\Delta,1}=\Phi_{\Delta,2}=\Phi_\Delta(t)$, integrating by parts the
cross terms gives
\begin{equation}
\begin{aligned}
  \mathcal{L}={}&\sum_{i=1,2}\Big[\frac{C}{2}\dot{\Phi}_i^2
  +\frac{C_p}{2}\big(\dot{\Phi}_i-V_p\big)^2\\
  &\qquad+C_{\mathrm{ext}}\,\Phi_i\ddot{\Phi}_\Delta-U_i(\phi_i,t)\Big]\\
  &+\frac{C_c}{2}\big(\dot{\Phi}_1-\dot{\Phi}_2\big)^2
  -\frac{(\Phi_1-\Phi_2)^2}{2L_c},
\end{aligned}
\label{eq:app_L}
\end{equation}
with $C=C_s+2C_J$ and $C_{\mathrm{ext}}=C_s+C_p$, up to total derivatives
and functions of time. The term $C_{\mathrm{ext}}\Phi_i\ddot{\Phi}_\Delta$
is the charge drive induced by the pumped $\Delta$-flux; it is symmetric
in $i$.

In configuration A both elements are biased identically,
$\phi_{\Sigma,i}=-\pi/2+\epsilon(t)$ and $\phi_{\Delta,i}=-\pi/2+\eta(t)$,
with $\epsilon(t)=\epsilon_p\sin(\omega_p t)$ and
$\eta(t)=\eta_p\sin(\omega_p t)$. Shifting $\vartheta_i=\phi_i+\pi/2$ and
absorbing the charge drive, each element contributes
$C\dot{\Phi}_i^2/2-U_i^{\mathrm{eff}}$ with
\begin{equation}
  U_i^{\mathrm{eff}}=\frac{E_L}{2}\vartheta_i^2-E_L^\eta\,\eta(t)\,\vartheta_i
  -2E_J\sin[\epsilon(t)]\sin\vartheta_i,
  \label{eq:app_Ueff}
\end{equation}
with $E_L^\eta=E_L(1-\omega_p^2/\omega_L^2)$ and $\omega_L^2=1/(LC_{\mathrm{ext}})$. The factor
$(1-\omega_p^2/\omega_L^2)$ is a small charge-drive correction
($\simeq0.98$ for typical parameters). The point $\vartheta_i=0$ is the
static equilibrium; at $\epsilon=0$ the junctions do not contribute to the
linear inductance. Equation~\eqref{eq:app_Ueff} is the classical,
single-element counterpart of the ATS potential of
Ref.~\cite{Carde2025Feb}.

\subsection{Normal modes and classical Hamiltonian}

We introduce the symmetric and antisymmetric coordinates
$\Psi_\pm=\phi_0(\vartheta_1\pm\vartheta_2)/2$. Using
$2\cos\theta_a\sin\theta_b=\sin(\theta_b+\theta_a)+\sin(\theta_b-\theta_a)$,
the Lagrangian separates into
\begin{equation}
\begin{aligned}
  \mathcal{L}={}&\frac{C_+}{2}\dot{\Psi}_+^2-\frac{\Psi_+^2}{2L_+}
  +\frac{C_-}{2}\dot{\Psi}_-^2-\frac{\Psi_-^2}{2L_-}\\
  &-2C_pV_p\dot{\Psi}_+
  +2E_L^\eta\,\eta(t)\frac{\Psi_+}{\phi_0}\\
  &+4E_J\sin[\epsilon(t)]\cos\frac{\Psi_-}{\phi_0}\sin\frac{\Psi_+}{\phi_0},
\end{aligned}
\label{eq:app_Lnormal}
\end{equation}
with
\begin{equation}
\begin{aligned}
  C_+&=2\bar{C}, & L_+&=\frac{L}{2},\\
  C_-&=2(\bar{C}+2C_c), & \frac{1}{L_-}&=\frac{2}{L}+\frac{4}{L_c},
\end{aligned}
\end{equation}
and $\bar{C}=C+C_p$. Mirror symmetry forbids any $\Psi_+\Psi_-$ or
$\dot{\Psi}_+\dot{\Psi}_-$ cross term, so $\Psi_\pm$ are exactly the
normal modes. Because the vector $(1,1)$ is an eigenvector of the
capacitance matrix and the $\Delta$-fluxes are identical, both the port
and the $\Delta$-pump couple only to $\Psi_+$.

With conjugate charges $Q_+=C_+\dot{\Psi}_+-2C_pV_p$ and
$Q_-=C_-\dot{\Psi}_-$ (i.e.\ $Q_\pm=Q_1\pm Q_2$), the Legendre transform
gives
\begin{equation}
\begin{aligned}
  \mathcal{H}={}&\frac{Q_+^2}{2C_+}+\frac{\Psi_+^2}{2L_+}
  +\frac{Q_-^2}{2C_-}+\frac{\Psi_-^2}{2L_-}\\
  &+\frac{C_p}{\bar{C}}V_p(t)\,Q_+
  -2E_L^\eta\,\eta(t)\frac{\Psi_+}{\phi_0}\\
  &-4E_J\sin[\epsilon(t)]\cos\frac{\Psi_-}{\phi_0}\sin\frac{\Psi_+}{\phi_0}.
\end{aligned}
\label{eq:app_H_bis}
\end{equation}
The linear frequencies are
\begin{equation}
\begin{aligned}
  \omega_b&=\frac{1}{\sqrt{L_+C_+}}=\frac{1}{\sqrt{L\bar{C}}},\\
  \omega_a&=\frac{1}{\sqrt{L_-C_-}}=\sqrt{\frac{1/L+2/L_c}{\bar{C}+2C_c}} .
\end{aligned}
\end{equation}
The antisymmetric mode $\Psi_-$ is the memory $a$ and the symmetric mode
$\Psi_+$ is the buffer $b$. Without $L_c$, the coupling capacitance
yields $\omega_a<\omega_b$; adding $L_c$ raises $\omega_a$.

\paragraph{Cancellation of the linear buffer drive.}
The nonlinear term contains the linear piece
$-4E_J\sin[\epsilon(t)]\Psi_+/\phi_0$. Using
$\sin\epsilon(t)=2\sum_{k>0,\,\mathrm{odd}}J_k(\epsilon_p)\sin(k\omega_pt)$,
the total linear drive on the buffer at $\omega_p$ is
$-[2E_L^\eta\eta_p+8E_JJ_1(\epsilon_p)]\sin(\omega_pt)\Psi_+/\phi_0$,
which vanishes for
\begin{equation}
  \eta_p=-\frac{4E_JJ_1(\epsilon_p)}{E_L^\eta}
  \;\xrightarrow{\;\epsilon_p\ll1\;}\;
  \frac{\eta_p}{\epsilon_p}=-\frac{2E_J}{E_L^\eta},
  \label{eq:app_eta}
\end{equation}
the two-element counterpart of the single-ATS
$g_{1,1}=0$ condition of Ref.~\cite{Carde2025Feb}. The residual linear drive is
the third harmonic, $-8E_JJ_3(\epsilon_p)\sin(3\omega_pt)\Psi_+/\phi_0$,
smaller by $J_3/J_1\simeq\epsilon_p^2/24$. This removes the circulating
pump photons in the buffer, and it does so with a parity-even drive.

\paragraph{Configuration B.}
Alternatively, the elements can be pumped in antiphase,
$\phi_{\Sigma,1,2}=-\pi/2\pm\epsilon(t)$, with the same static bias. The
nonlinearity then becomes
$-4E_J\sin\epsilon\,\cos(\Psi_+/\phi_0)\sin(\Psi_-/\phi_0)$, so the roles
of the modes are exchanged: the memory must be $\Psi_+$ and the buffer
$\Psi_-$, and the port must couple to $Q_-$ only (e.g.\ through a balun).
Because of this port requirement, configuration A is preferred.

\subsection{Quantization}
\label{app:quantization}

We promote $[\hat{\Psi}_\pm,\hat{Q}_\pm]=i\hbar$ and write the line
voltage as $V_p(t)=V_d\cos(\omega_dt)+\hat{V}_B$, a classical buffer drive
plus the transmission-line bath operator. With
$\hat{\Psi}_-=\Psi_a^{\mathrm{zpf}}\hat{x}_a$,
$\hat{Q}_-=Q_a^{\mathrm{zpf}}\hat{y}_a$ (and likewise for $b$ with
$\Psi_+,Q_+$), $\hat{x}_a=\hat{a}+\hat{a}^\dagger$,
$\hat{y}_a=-i(\hat{a}-\hat{a}^\dagger)$ (same for $b$),
$\Psi^{\mathrm{zpf}}=\sqrt{\hbar Z/2}$, $Q^{\mathrm{zpf}}=\sqrt{\hbar/2Z}$
and $Z_{a,b}=\sqrt{L_\mp/C_\mp}$, the phase zero-point fluctuations are
\begin{equation}
  \varphi_a=\left(\frac{2E_{C_-}}{E_{L_-}}\right)^{1/4},\qquad
  \varphi_b=\left(\frac{2E_{C_+}}{E_{L_+}}\right)^{1/4},
\end{equation}
with $E_{C_\pm}=e^2/2C_\pm$ and $E_{L_\pm}=\phi_0^2/L_\pm$. The island
phases are $\hat{\vartheta}_{1,2}=\varphi_b\hat{x}_b\pm\varphi_a\hat{x}_a$:
the memory enters the two elements with opposite signs and the buffer
with the same sign. Dropping zero-point energies, the system and
system--bath Hamiltonians read
\begin{equation}
\begin{aligned}
  \frac{\hat{H}_s}{\hbar}={}&\omega_a\hat{a}^\dagger\hat{a}
  +\omega_b\hat{b}^\dagger\hat{b}\\
  &+2\epsilon_d\cos(\omega_dt)\,\hat{y}_b
  -2\frac{E_L^\eta}{\hbar}\eta(t)\,\varphi_b\hat{x}_b\\
  &-4\frac{E_J}{\hbar}\sin[\epsilon(t)]
  \cos(\varphi_a\hat{x}_a)\sin(\varphi_b\hat{x}_b),\\
  \frac{\hat{H}_{sB}}{\hbar}={}&\hat{y}_b\otimes\hat{B},
\end{aligned}
\label{eq:app_Hq}
\end{equation}
with $\epsilon_d=C_pQ_b^{\mathrm{zpf}}V_d/(2\hbar\bar{C})$ and
$\hat{B}=C_pQ_b^{\mathrm{zpf}}\hat{V}_B/(\hbar\bar{C})$.

Equation~\eqref{eq:app_Hq} has the same structure as the single-ATS
Hamiltonian of Ref.~\cite{Carde2025Feb}, with three differences:
(i)~the nonlinearity is the target form~\eqref{eq:app_target}, even under
$\hat{x}_a\to-\hat{x}_a$ at all times;
(ii)~the memory does not appear in either the drive or the bath coupling
($u=0$ exactly), i.e.\ it is a dark mode of the port;
(iii)~there is no linear drive on $\hat{x}_a$ for any $\eta_p$, so the
memory counterpart of $g_{1,1}$ vanishes identically. The linear pump
terms act on the buffer only and cancel at the first harmonic for the
choice~\eqref{eq:app_eta}. Two parity-protected operating points are
therefore available: $\eta_p=0$, the counterpart of the baseline
single-ATS operation, and $\eta_p$ given by Eq.~\eqref{eq:app_eta}, the
counterpart of its mitigated working point, which in addition has much
smaller Stark and Kerr shifts.

\section{Transient chaos}
\label{sec:transient_chaos}


We integrate the classical mean-field equations of the memory and buffer amplitudes $a(t)$ and $b(t)$. The initial conditions form a uniform $10\times10$ grid in the memory phase space.

Each trajectory is sampled once per memory period $T_a = 2\pi/\omega_a$,
this is the Poincaré map of the memory at period $T_a$. A periodic orbit of period $T_a$ maps to a fixed point, a period-doubled orbit to two points, and a chaotic orbit to an extended set.

Each sampled trajectory is cut into $K = \lfloor N_s/M \rfloor$ consecutive windows of $M$ points,
\begin{equation}
    \mathcal{C}^{(j)}_k = \bigl\{ a^{(j)}_n \,:\, (k-1)M \le n < kM \bigr\}, \qquad k = 1,\dots,K .
\end{equation}

We will compute the entropy of these trajectory cuts and we start by cutting phase space in cells $\mathcal{B}_{lm}$. The same partition is used for every window, so that entropies can be compared. For each trajectory and window we compute the occupation probabilities and the Shannon entropy
\begin{equation}
\begin{split}
    p^{(j,k)}_{lm} =&\; \frac{1}{M}\,\#\bigl\{ a \in \mathcal{C}^{(j)}_k \,:\, a \in \mathcal{B}_{lm} \bigr\},
    \\
    H_{j,k} =&\; -\sum_{l,m} p^{(j,k)}_{lm} \ln p^{(j,k)}_{lm},
\end{split}
\end{equation}
with $0\ln 0 = 0$. By construction $0 \le H_{j,k} \le \ln M \simeq 5.30$.
\begin{itemize}
    \item A window whose $M$ points all fall in the same cell, such as a period-$T_a$ orbit, gives $H = 0$.
    \item A period-$2T_a$ orbit gives $H = \ln 2$.
    \item A window whose points spread over $M$ distinct cells, as for a chaotic trajectory, saturates at $H \simeq \ln M$.
\end{itemize}

\paragraph*{Classification and fractions.}
A trajectory is counted as chaotic in window $k$ if $H_{j,k} \ge H_\text{c} > \ln 2$, and as regular if $H_{j,k} \le H_\text{c}$. The regular criterion selects attractors visited by at most two cells, i.e.\ period-$T_a$ and period-$2T_a$ orbits. The plotted quantities are the fractions of the ensemble in each class,
\begin{equation}
    f_\text{c}(t_k) = \frac{1}{N}\,\#\bigl\{ j : H_{j,k} \ge H_\text{c} \bigr\},
    \qquad
    f_\text{r}(t_k) = \frac{1}{N}\,\#\bigl\{ j : H_{j,k} \le H_\text{r} \bigr\},
\end{equation}
shown as solid lines against $\kappa_b t_k$.

\TMD{Transient chaos literature, reference for exponential decay, Difficulty for the computation of Lyapunov exponents, work already on quantum transient chaos? Maybe the chaotic chaos is never transient?}

\begin{figure}[tbp]
    \centering
    \includegraphics[width=\columnwidth]{Chaos/figures_article/transient_chaos_12_10_2026.pdf}
    \caption{Decay of the proportion of chaotic trajectories of the system. The solid lines are exponential fits.}
    \label{fig:transient_chaos}
\end{figure}


\bibliography{biblio}

\end{document}